\chapter{Second quantization}
\label{chap:secondquant}

Chapter~\ref{chap:mbphys} ended with a warning.  Every expectation value there
was obtained by writing out the antisymmetriser, sorting through $N!$
permutations, and arguing term by term about which overlaps vanish.  For two
particles this is an exercise; for the eight-operator strings of
Section~\ref{sec:expval} it is already unpleasant; for the excited
determinants that perturbation theory and coupled-cluster theory require it is
impossible.  Exercise~2 of the week-36 session made the point sharply: the
same matrix elements fall out of the anticommutation relations in three lines.

This chapter develops that observation into a formalism.  Antisymmetry, which
was an explicit and increasingly costly bookkeeping task in
Chapter~\ref{chap:mbphys}, is built once and for all into the algebra of the
creation and annihilation operators.  The Slater determinant of
Eq.~(\ref{eq:SlatDet}) becomes a string of creation operators acting on the
vacuum, the Condon-Slater rules of Section~\ref{sec:particlehole} become
consequences of Wick's theorem, and calculations that were barely feasible by
hand for $N=3$ become mechanical for arbitrary $N$.

We introduce the time-independent  operators
$a_\alpha^{\dagger}$ and $a_\alpha$   which create and annihilate, respectively, a particle 
in the single-particle state 
$\varphi_\alpha$. 
We define the fermion creation operator
$a_\alpha^{\dagger}$ 
\begin{equation}
	a_\alpha^{\dagger}|0\rangle \equiv  |\alpha\rangle  \label{eq:2-1a},
\end{equation}
and
\begin{equation}
	a_\alpha^{\dagger}|\alpha_1\dots \alpha_n\rangle_{\mathrm{AS}} \equiv  |\alpha\alpha_1\dots \alpha_n\rangle_{\mathrm{AS}} \label{eq:2-1b}
\end{equation}

In Eq.~(\ref{eq:2-1a}) 
the operator  $a_\alpha^{\dagger}$  acts on the vacuum state 
$|0\rangle$, which does not contain any particles. Alternatively, we could define  a closed-shell nucleus or atom as our new vacuum, but then
we need to introduce the particle-hole  formalism, see the discussion to come. 

In Eq.~(\ref{eq:2-1b}) $a_\alpha^{\dagger}$ acts on an antisymmetric $n$-particle state and 
creates an antisymmetric $(n+1)$-particle state, where the one-body state 
$\varphi_\alpha$ is occupied, under the condition that
$\alpha \ne \alpha_1, \alpha_2, \dots, \alpha_n$. 
It follows that we can express an antisymmetric state as the product of the creation
operators acting on the vacuum state.  
\begin{equation}
	|\alpha_1\dots \alpha_n\rangle_{\mathrm{AS}} = a_{\alpha_1}^{\dagger} a_{\alpha_2}^{\dagger} \dots a_{\alpha_n}^{\dagger} |0\rangle \label{eq:2-2}
\end{equation}

It is easy to derive the commutation and anticommutation rules  for the fermionic creation operators 
$a_\alpha^{\dagger}$. Using the antisymmetry of the states 
(\ref{eq:2-2})
\begin{equation}
	|\alpha_1\dots \alpha_i\dots \alpha_k\dots \alpha_n\rangle_{\mathrm{AS}} = 
		- |\alpha_1\dots \alpha_k\dots \alpha_i\dots \alpha_n\rangle_{\mathrm{AS}} \label{eq:2-3a}
\end{equation}
we obtain
\begin{equation}
	 a_{\alpha_i}^{\dagger}  a_{\alpha_k}^{\dagger} = - a_{\alpha_k}^{\dagger} a_{\alpha_i}^{\dagger} \label{eq:2-3b}
\end{equation}

Using the Pauli principle
\begin{equation}
	|\alpha_1\dots \alpha_i\dots \alpha_i\dots \alpha_n\rangle_{\mathrm{AS}} = 0 \label{eq:2-4a}
\end{equation}
it follows that
\begin{equation}
	a_{\alpha_i}^{\dagger}  a_{\alpha_i}^{\dagger} = 0. \label{eq:2-4b}
\end{equation}
If we combine Eqs.~(\ref{eq:2-3b}) and (\ref{eq:2-4b}), we obtain the well-known anti-commutation rule
\begin{equation}
	a_{\alpha}^{\dagger}  a_{\beta}^{\dagger} + a_{\beta}^{\dagger}  a_{\alpha}^{\dagger} \equiv 
		\{a_{\alpha}^{\dagger},a_{\beta}^{\dagger}\} = 0 \label{eq:2-5}
\end{equation}

The hermitian conjugate  of $a_\alpha^{\dagger}$ is
\begin{equation}
	a_{\alpha} = ( a_{\alpha}^{\dagger} )^{\dagger} \label{eq:2-6}
\end{equation}
If we take the hermitian conjugate of Eq.~(\ref{eq:2-5}), we arrive at 
\begin{equation}
	\{a_{\alpha},a_{\beta}\} = 0 \label{eq:2-7}
\end{equation}

What is the physical interpretation of the operator $a_\alpha$ and what is the effect of 
$a_\alpha$ on a given state $|\alpha_1\alpha_2\dots\alpha_n\rangle_{\mathrm{AS}}$? 
Consider the following matrix element
\begin{equation}
	\langle\alpha_1\alpha_2 \dots \alpha_n|a_\alpha|\alpha_1'\alpha_2' \dots \alpha_m'\rangle \label{eq:2-8}
\end{equation}
where both sides are antisymmetric. We  distinguish between two cases. The first (1) is when
$\alpha \in \{\alpha_i\}$. Using the Pauli principle of Eq.~(\ref{eq:2-4a}) it follows
\begin{equation}
		\langle\alpha_1\alpha_2 \dots \alpha_n|a_\alpha = 0 \label{eq:2-9a}
\end{equation}
The second (2) case is when $\alpha \notin \{\alpha_i\}$. It follows that an hermitian conjugation
\begin{equation}
		\langle \alpha_1\alpha_2 \dots \alpha_n|a_\alpha = \langle\alpha\alpha_1\alpha_2 \dots \alpha_n|  \label{eq:2-9b}
\end{equation}

Eq.~(\ref{eq:2-9b}) holds for case (1) since the lefthand side is zero due to the Pauli principle. We write
Eq.~(\ref{eq:2-8}) as
\begin{equation}
	\langle\alpha_1\alpha_2 \dots \alpha_n|a_\alpha|\alpha_1'\alpha_2' \dots \alpha_m'\rangle = 
	\langle \alpha_1\alpha_2 \dots \alpha_n|\alpha\alpha_1'\alpha_2' \dots \alpha_m'\rangle \label{eq:2-10}
\end{equation}
Here we must have $m = n+1$ if Eq.~(\ref{eq:2-10}) has to be trivially different from zero.

For the last case, the minus and plus signs apply when the sequence 
$\alpha ,\alpha_1, \alpha_2, \dots, \alpha_n$ and 
$\alpha_1', \alpha_2', \dots, \alpha_{n+1}'$ are related to each other via even and odd permutations.
If we assume that  $\alpha \notin \{\alpha_i\}$ we obtain 
\begin{equation}
	\langle\alpha_1\alpha_2 \dots \alpha_n|a_\alpha|\alpha_1'\alpha_2' \dots \alpha_{n+1}'\rangle = 0 \label{eq:2-12}
\end{equation}
when $\alpha \in \{\alpha_i'\}$. If $\alpha \notin \{\alpha_i'\}$, we obtain
\begin{equation}
	a_\alpha\underbrace{|\alpha_1'\alpha_2' \dots \alpha_{n+1}'}\rangle_{\neq \alpha} = 0 \label{eq:2-13a}
\end{equation}
and in particular
\begin{equation}
	a_\alpha |0\rangle = 0 \label{eq:2-13b}
\end{equation}

If $\{\alpha\alpha_i\} = \{\alpha_i'\}$, performing the right permutations, the sequence
$\alpha ,\alpha_1, \alpha_2, \dots, \alpha_n$ is identical with the sequence
$\alpha_1', \alpha_2', \dots, \alpha_{n+1}'$. This results in
\begin{equation}
	\langle\alpha_1\alpha_2 \dots \alpha_n|a_\alpha|\alpha\alpha_1\alpha_2 \dots \alpha_{n}\rangle = 1 \label{eq:2-14}
\end{equation}
and thus
\begin{equation}
	a_\alpha |\alpha\alpha_1\alpha_2 \dots \alpha_{n}\rangle = |\alpha_1\alpha_2 \dots \alpha_{n}\rangle \label{eq:2-15}
\end{equation}

The action of the operator 
$a_\alpha$ from the left on a state vector  is to to remove  one particle in the state
$\alpha$. 
If the state vector does not contain the single-particle state $\alpha$, the outcome of the operation is zero.
The operator  $a_\alpha$ is normally called for a destruction or annihilation operator.

The next step is to establish the  commutator algebra of $a_\alpha^{\dagger}$ and
$a_\beta$. 

The action of the anti-commutator 
$\{a_\alpha^{\dagger}$,$a_\alpha\}$ on a given $n$-particle state is
\begin{align}
	a_\alpha^{\dagger} a_\alpha \underbrace{|\alpha_1\alpha_2 \dots \alpha_{n}\rangle}_{\neq \alpha} &= 0 \nonumber \\
	a_\alpha a_\alpha^{\dagger} \underbrace{|\alpha_1\alpha_2 \dots \alpha_{n}\rangle}_{\neq \alpha} &=
	a_\alpha \underbrace{|\alpha \alpha_1\alpha_2 \dots \alpha_{n}\rangle}_{\neq \alpha} = 
	\underbrace{|\alpha_1\alpha_2 \dots \alpha_{n}\rangle}_{\neq \alpha} \label{eq:2-16a}
\end{align}
if the single-particle state $\alpha$ is not contained in the state.

 If it is present
we arrive at
\begin{align}
	a_\alpha^{\dagger} a_\alpha |\alpha_1\alpha_2 \dots \alpha_{k}\alpha \alpha_{k+1} \dots \alpha_{n-1}\rangle &=
	a_\alpha^{\dagger} a_\alpha (-1)^k |\alpha \alpha_1\alpha_2 \dots \alpha_{n-1}\rangle \nonumber \\
	= (-1)^k |\alpha \alpha_1\alpha_2 \dots \alpha_{n-1}\rangle &=
	|\alpha_1\alpha_2 \dots \alpha_{k}\alpha \alpha_{k+1} \dots \alpha_{n-1}\rangle \nonumber \\
	a_\alpha a_\alpha^{\dagger}|\alpha_1\alpha_2 \dots \alpha_{k}\alpha \alpha_{k+1} \dots \alpha_{n-1}\rangle &= 0 \label{eq:2-16b}
\end{align}
From Eqs.~(\ref{eq:2-16a}) and  (\ref{eq:2-16b}) we arrive at 
\begin{equation}
	\{a_\alpha^{\dagger} , a_\alpha \} = a_\alpha^{\dagger} a_\alpha + a_\alpha a_\alpha^{\dagger} = 1 \label{eq:2-17}
\end{equation}

The action of $\left\{a_\alpha^{\dagger}, a_\beta\right\}$, with 
$\alpha \ne \beta$ on a given state yields three possibilities. 
The first case is a state vector which contains both $\alpha$ and $\beta$, then either 
$\alpha$ or $\beta$ and finally none of them.

The first case results in
\begin{align}
	a_\alpha^{\dagger} a_\beta |\alpha\beta\alpha_1\alpha_2 \dots \alpha_{n-2}\rangle = 0 \nonumber \\
	a_\beta a_\alpha^{\dagger} |\alpha\beta\alpha_1\alpha_2 \dots \alpha_{n-2}\rangle = 0 \label{eq:2-18a}
\end{align}
while the second case gives
\begin{align}
	 a_\alpha^{\dagger} a_\beta |\beta \underbrace{\alpha_1\alpha_2 \dots \alpha_{n-1}}_{\neq \alpha}\rangle =& 
	 	|\alpha \underbrace{\alpha_1\alpha_2 \dots \alpha_{n-1}}_{\neq  \alpha}\rangle \nonumber \\
	a_\beta a_\alpha^{\dagger} |\beta \underbrace{\alpha_1\alpha_2 \dots \alpha_{n-1}}_{\neq \alpha}\rangle =&
		a_\beta |\alpha\beta\underbrace{\beta \alpha_1\alpha_2 \dots \alpha_{n-1}}_{\neq \alpha}\rangle \nonumber \\
	=& - |\alpha\underbrace{\alpha_1\alpha_2 \dots \alpha_{n-1}}_{\neq \alpha}\rangle \label{eq:2-18b}
\end{align}

Finally if the state vector does not contain $\alpha$ and $\beta$
\begin{align}
	a_\alpha^{\dagger} a_\beta |\underbrace{\alpha_1\alpha_2 \dots \alpha_{n}}_{\neq \alpha,\beta}\rangle &=& 0 \nonumber \\
	a_\beta a_\alpha^{\dagger} |\underbrace{\alpha_1\alpha_2 \dots \alpha_{n}}_{\neq \alpha,\beta}\rangle &=& 
		a_\beta |\alpha \underbrace{\alpha_1\alpha_2 \dots \alpha_{n}}_{\neq \alpha,\beta}\rangle = 0 \label{eq:2-18c}
\end{align}
For all three cases we have
\begin{equation}
	\{a_\alpha^{\dagger},a_\beta \} = a_\alpha^{\dagger} a_\beta + a_\beta a_\alpha^{\dagger} = 0, \quad \alpha \neq \beta \label{eq:2-19}
\end{equation}

We can summarize  our findings in Eqs.~(\ref{eq:2-17}) and (\ref{eq:2-19}) as 
\begin{equation}
	\{a_\alpha^{\dagger},a_\beta \} = \delta_{\alpha\beta} \label{eq:2-20}
\end{equation}
with  $\delta_{\alpha\beta}$ is the Kroenecker $\delta$-symbol.

The properties of the creation and annihilation operators can be summarized as (for fermions)
\[
	a_\alpha^{\dagger}|0\rangle \equiv  |\alpha\rangle,
\]
and
\[
	a_\alpha^{\dagger}|\alpha_1\dots \alpha_n\rangle_{\mathrm{AS}} \equiv  |\alpha\alpha_1\dots \alpha_n\rangle_{\mathrm{AS}}. 
\]
from which follows
\[
        |\alpha_1\dots \alpha_n\rangle_{\mathrm{AS}} = a_{\alpha_1}^{\dagger} a_{\alpha_2}^{\dagger} \dots a_{\alpha_n}^{\dagger} |0\rangle.
\]

The hermitian conjugate has the folowing properties
\[
        a_{\alpha} = ( a_{\alpha}^{\dagger} )^{\dagger}.
\]
Finally we found 
\[
	a_\alpha\underbrace{|\alpha_1'\alpha_2' \dots \alpha_{n+1}'}\rangle_{\neq \alpha} = 0, \quad
		\textrm{in particular } a_\alpha |0\rangle = 0,
\]
and
\[
 a_\alpha |\alpha\alpha_1\alpha_2 \dots \alpha_{n}\rangle = |\alpha_1\alpha_2 \dots \alpha_{n}\rangle,
\]
and the corresponding commutator algebra
\[
	\{a_{\alpha}^{\dagger},a_{\beta}^{\dagger}\} = \{a_{\alpha},a_{\beta}\} = 0 \hspace{0.5cm} 
\{a_\alpha^{\dagger},a_\beta \} = \delta_{\alpha\beta}.
\]

\section{One-body operators in second quantization}

The result we are about to derive was obtained once already, in
Section~\ref{sec:expval}, by integrating over permutations of the Hartree
function.  It is worth comparing the two routes as we go: what took a page of
careful argument there follows here from the anticommutation relations alone.


A very useful operator is the so-called number-operator.  Most physics cases  we will
study in this text conserve the total number of particles.  The number operator is therefore
a useful quantity which allows us to test that our many-body formalism  conserves the number of particles.
In for example $(d,p)$ or $(p,d)$ reactions it is important to be able to describe quantum mechanical states
where particles get added or removed.
A creation operator $a_\alpha^{\dagger}$ adds one particle to the single-particle state
$\alpha$ of a give many-body state vector, while an annihilation operator $a_\alpha$ 
removes a particle from a single-particle
state $\alpha$. 

Let us consider an operator proportional with $a_\alpha^{\dagger} a_\beta$ and 
$\alpha=\beta$. It acts on an $n$-particle state 
resulting in
\begin{equation}
	a_\alpha^{\dagger} a_\alpha |\alpha_1\alpha_2 \dots \alpha_{n}\rangle = 
	\begin{cases}
		0  &\alpha \notin \{\alpha_i\} \\
		\\
		|\alpha_1\alpha_2 \dots \alpha_{n}\rangle & \alpha \in \{\alpha_i\}
	\end{cases}
\end{equation}
Summing over all possible one-particle states we arrive at
\begin{equation}
	\left( \sum_\alpha a_\alpha^{\dagger} a_\alpha \right) |\alpha_1\alpha_2 \dots \alpha_{n}\rangle = 
	n |\alpha_1\alpha_2 \dots \alpha_{n}\rangle \label{eq:2-21}
\end{equation}

The operator 
\begin{equation}
	\hat{N} = \sum_\alpha a_\alpha^{\dagger} a_\alpha \label{eq:2-22}
\end{equation}
is called the number operator since it counts the number of particles in a give state vector when it acts 
on the different single-particle states.  It acts on one single-particle state at the time and falls 
therefore under category one-body operators.
Next we look at another important one-body operator, namely $\hat{H}_0$ and study its operator form in the 
occupation number representation.

We want to obtain an expression for a one-body operator which conserves the number of particles.
Here we study the one-body operator for the kinetic energy plus an eventual external one-body potential.
The action of this operator on a particular $n$-body state with its pertinent expectation value has already
been studied in coordinate  space.
In coordinate space the operator reads
\begin{equation}
	\hat{H}_0 = \sum_i \hat{h}_0(x_i) \label{eq:2-23}
\end{equation}
and the anti-symmetric $n$-particle Slater determinant is defined as 
\[
\Phi(x_1, x_2,\dots ,x_n,\alpha_1,\alpha_2,\dots, \alpha_n)= \frac{1}{\sqrt{n!}} \sum_p (-1)^p\hat{P}\psi_{\alpha_1}(x_1)\psi_{\alpha_2}(x_2) \dots \psi_{\alpha_n}(x_n).
\]

Defining
\begin{equation}
	\hat{h}_0(x_i) \psi_{\alpha_i}(x_i) = \sum_{\alpha_k'} \psi_{\alpha_k'}(x_i) \langle\alpha_k'|\hat{h}_0|\alpha_k\rangle \label{eq:2-25}
\end{equation}
we can easily  evaluate the action of $\hat{H}_0$ on each product of one-particle functions in Slater determinant.
From Eq.~(\ref{eq:2-25})  we obtain the following result without  permuting any particle pair 
\begin{align}
	&& \left( \sum_i \hat{h}_0(x_i) \right) \psi_{\alpha_1}(x_1)\psi_{\alpha_2}(x_2) \dots \psi_{\alpha_n}(x_n) \nonumber \\
	& =&\sum_{\alpha_1'} \langle \alpha_1'|\hat{h}_0|\alpha_1\rangle 
		\psi_{\alpha_1'}(x_1)\psi_{\alpha_2}(x_2) \dots \psi_{\alpha_n}(x_n) \nonumber \\
	&+&\sum_{\alpha_2'} \langle \alpha_2'|\hat{h}_0|\alpha_2\rangle
		\psi_{\alpha_1}(x_1)\psi_{\alpha_2'}(x_2) \dots \psi_{\alpha_n}(x_n) \nonumber \\
	&+& \dots \nonumber \\
	&+&\sum_{\alpha_n'} \langle \alpha_n'|\hat{h}_0|\alpha_n\rangle
		\psi_{\alpha_1}(x_1)\psi_{\alpha_2}(x_2) \dots \psi_{\alpha_n'}(x_n) \label{eq:2-26}
\end{align}

If we interchange particles $1$ and $2$  we obtain
\begin{align}
	&& \left( \sum_i \hat{h}_0(x_i) \right) \psi_{\alpha_1}(x_2)\psi_{\alpha_1}(x_2) \dots \psi_{\alpha_n}(x_n) \nonumber \\
	& =&\sum_{\alpha_2'} \langle \alpha_2'|\hat{h}_0|\alpha_2\rangle 
		\psi_{\alpha_1}(x_2)\psi_{\alpha_2'}(x_1) \dots \psi_{\alpha_n}(x_n) \nonumber \\
	&+&\sum_{\alpha_1'} \langle \alpha_1'|\hat{h}_0|\alpha_1\rangle
		\psi_{\alpha_1'}(x_2)\psi_{\alpha_2}(x_1) \dots \psi_{\alpha_n}(x_n) \nonumber \\
	&+& \dots \nonumber \\
	&+&\sum_{\alpha_n'} \langle \alpha_n'|\hat{h}_0|\alpha_n\rangle
		\psi_{\alpha_1}(x_2)\psi_{\alpha_1}(x_2) \dots \psi_{\alpha_n'}(x_n) \label{eq:2-27}
\end{align}

We can continue by computing all possible permutations. 
We rewrite also our Slater determinant in its second quantized form and skip the dependence on the quantum numbers $x_i.$
Summing up all contributions and taking care of all phases
$(-1)^p$ we arrive at 
\begin{align}
	\hat{H}_0|\alpha_1,\alpha_2,\dots, \alpha_n\rangle &=& \sum_{\alpha_1'}\langle \alpha_1'|\hat{h}_0|\alpha_1\rangle
		|\alpha_1'\alpha_2 \dots \alpha_{n}\rangle \nonumber \\
	&+& \sum_{\alpha_2'} \langle \alpha_2'|\hat{h}_0|\alpha_2\rangle
		|\alpha_1\alpha_2' \dots \alpha_{n}\rangle \nonumber \\
	&+& \dots \nonumber \\
	&+& \sum_{\alpha_n'} \langle \alpha_n'|\hat{h}_0|\alpha_n\rangle
		|\alpha_1\alpha_2 \dots \alpha_{n}'\rangle \label{eq:2-28}
\end{align}

In Eq.~(\ref{eq:2-28}) 
we have expressed the action of the one-body operator
of Eq.~(\ref{eq:2-23}) on the  $n$-body state in its second quantized form.
This equation can be further manipulated if we use the properties of the creation and annihilation operator
on each primed quantum number, that is
\begin{equation}
	|\alpha_1\alpha_2 \dots \alpha_k' \dots \alpha_{n}\rangle = 
		a_{\alpha_k'}^{\dagger}  a_{\alpha_k} |\alpha_1\alpha_2 \dots \alpha_k \dots \alpha_{n}\rangle \label{eq:2-29}
\end{equation}
Inserting this in the right-hand side of Eq.~(\ref{eq:2-28}) results in
\begin{align}
	\hat{H}_0|\alpha_1\alpha_2 \dots \alpha_{n}\rangle &=& \sum_{\alpha_1'}\langle \alpha_1'|\hat{h}_0|\alpha_1\rangle
		a_{\alpha_1'}^{\dagger}  a_{\alpha_1} |\alpha_1\alpha_2 \dots \alpha_{n}\rangle \nonumber \\
	&+& \sum_{\alpha_2'} \langle \alpha_2'|\hat{h}_0|\alpha_2\rangle
		a_{\alpha_2'}^{\dagger}  a_{\alpha_2} |\alpha_1\alpha_2 \dots \alpha_{n}\rangle \nonumber \\
	&+& \dots \nonumber \\
	&+& \sum_{\alpha_n'} \langle \alpha_n'|\hat{h}_0|\alpha_n\rangle
		a_{\alpha_n'}^{\dagger}  a_{\alpha_n} |\alpha_1\alpha_2 \dots \alpha_{n}\rangle \nonumber \\
	&=& \sum_{\alpha, \beta} \langle \alpha|\hat{h}_0|\beta\rangle a_\alpha^{\dagger} a_\beta 
		|\alpha_1\alpha_2 \dots \alpha_{n}\rangle \label{eq:2-30a}
\end{align}

In the number occupation representation or second quantization we get the following expression for a one-body 
operator which conserves the number of particles
\begin{equation}
	\hat{H}_0 = \sum_{\alpha\beta} \langle \alpha|\hat{h}_0|\beta\rangle a_\alpha^{\dagger} a_\beta \label{eq:2-30b}
\end{equation}
Obviously, $\hat{H}_0$ can be replaced by any other one-body  operator which preserved the number
of particles. The stucture of the operator is therefore not limited to say the kinetic or single-particle energy only.

The opearator $\hat{H}_0$ takes a particle from the single-particle state $\beta$  to the single-particle state $\alpha$ 
with a probability for the transition given by the expectation value $\langle \alpha|\hat{h}_0|\beta\rangle$.

It is instructive to verify Eq.~(\ref{eq:2-30b}) by computing the expectation value of $\hat{H}_0$ 
between two single-particle states
\begin{equation}
	\langle \alpha_1|\hat{h}_0|\alpha_2\rangle = \sum_{\alpha\beta} \langle \alpha|\hat{h}_0|\beta\rangle
		\langle 0|a_{\alpha_1}a_\alpha^{\dagger} a_\beta a_{\alpha_2}^{\dagger}|0\rangle \label{eq:2-30c}
\end{equation}

Using the commutation relations for the creation and annihilation operators we have 
\begin{equation}
a_{\alpha_1}a_\alpha^{\dagger} a_\beta a_{\alpha_2}^{\dagger} = (\delta_{\alpha \alpha_1} - a_\alpha^{\dagger} a_{\alpha_1} )(\delta_{\beta \alpha_2} - a_{\alpha_2}^{\dagger} a_{\beta} ), \label{eq:2-30d}
\end{equation}
which results in
\begin{equation}
\langle 0|a_{\alpha_1}a_\alpha^{\dagger} a_\beta a_{\alpha_2}^{\dagger}|0\rangle = \delta_{\alpha \alpha_1} \delta_{\beta \alpha_2} \label{eq:2-30e}
\end{equation}
and
\begin{equation}
\langle \alpha_1|\hat{h}_0|\alpha_2\rangle = \sum_{\alpha\beta} \langle \alpha|\hat{h}_0|\beta\rangle\delta_{\alpha \alpha_1} \delta_{\beta \alpha_2} = \langle \alpha_1|\hat{h}_0|\alpha_2\rangle \label{eq:2-30f}
\end{equation}

\section{Two-body operators in second quantization}

The antisymmetrised matrix elements $\langle pq|\hat{v}|rs\rangle_{AS}$ that
appear below are exactly those of Eq.~(\ref{eq:2-antisymmatrix}), and the
$n^4$ of them are the storage problem discussed in
Section~\ref{sec:twobodycost}.  Nothing in this chapter changes that count;
what changes is the ease with which they are contracted.


Let us now derive the expression for our two-body interaction part, which also conserves the number of particles.
We can proceed in exactly the same way as for the one-body operator. In the coordinate representation our
two-body interaction part takes the following expression
\begin{equation}
	\hat{H}_I = \sum_{i < j} V(x_i,x_j) \label{eq:2-31}
\end{equation}
where the summation runs over distinct pairs. The term $V$ can be an interaction model for the nucleon-nucleon interaction
or the interaction between two electrons. It can also include additional two-body interaction terms. 

The action of this operator on a product of 
two single-particle functions is defined as 
\begin{equation}
	V(x_i,x_j) \psi_{\alpha_k}(x_i) \psi_{\alpha_l}(x_j) = \sum_{\alpha_k'\alpha_l'} 
		\psi_{\alpha_k}'(x_i)\psi_{\alpha_l}'(x_j) 
		\langle \alpha_k'\alpha_l'|\hat{v}|\alpha_k\alpha_l\rangle \label{eq:2-32}
\end{equation}

We can now let $\hat{H}_I$ act on all terms in the linear combination for $|\alpha_1\alpha_2\dots\alpha_n\rangle$. Without any permutations we have
\begin{align}
	&& \left( \sum_{i < j} V(x_i,x_j) \right) \psi_{\alpha_1}(x_1)\psi_{\alpha_2}(x_2)\dots \psi_{\alpha_n}(x_n) \nonumber \\
	&=& \sum_{\alpha_1'\alpha_2'} \langle \alpha_1'\alpha_2'|\hat{v}|\alpha_1\alpha_2\rangle
		\psi_{\alpha_1}'(x_1)\psi_{\alpha_2}'(x_2)\dots \psi_{\alpha_n}(x_n) \nonumber \\
	& +& \dots \nonumber \\
	&+& \sum_{\alpha_1'\alpha_n'} \langle \alpha_1'\alpha_n'|\hat{v}|\alpha_1\alpha_n\rangle
		\psi_{\alpha_1}'(x_1)\psi_{\alpha_2}(x_2)\dots \psi_{\alpha_n}'(x_n) \nonumber \\
	& +& \dots \nonumber \\
	&+& \sum_{\alpha_2'\alpha_n'} \langle \alpha_2'\alpha_n'|\hat{v}|\alpha_2\alpha_n\rangle
		\psi_{\alpha_1}(x_1)\psi_{\alpha_2}'(x_2)\dots \psi_{\alpha_n}'(x_n) \nonumber \\
	 & +& \dots \label{eq:2-33}
\end{align}
where on the rhs we have a term for each distinct pairs. 

For the other terms on the rhs we obtain similar expressions  and summing over all terms we obtain
\begin{align}
	H_I |\alpha_1\alpha_2\dots\alpha_n\rangle &=& \sum_{\alpha_1', \alpha_2'} \langle \alpha_1'\alpha_2'|\hat{v}|\alpha_1\alpha_2\rangle
		|\alpha_1'\alpha_2'\dots\alpha_n\rangle \nonumber \\
	&+& \dots \nonumber \\
	&+& \sum_{\alpha_1', \alpha_n'} \langle \alpha_1'\alpha_n'|\hat{v}|\alpha_1\alpha_n\rangle
		|\alpha_1'\alpha_2\dots\alpha_n'\rangle \nonumber \\
	&+& \dots \nonumber \\
	&+& \sum_{\alpha_2', \alpha_n'} \langle \alpha_2'\alpha_n'|\hat{v}|\alpha_2\alpha_n\rangle
		|\alpha_1\alpha_2'\dots\alpha_n'\rangle \nonumber \\
	 &+& \dots \label{eq:2-34}
\end{align}

We introduce second quantization via the relation
\begin{align}
	&& a_{\alpha_k'}^{\dagger} a_{\alpha_l'}^{\dagger} a_{\alpha_l} a_{\alpha_k} 
		|\alpha_1\alpha_2\dots\alpha_k\dots\alpha_l\dots\alpha_n\rangle \nonumber \\
	&=& (-1)^{k-1} (-1)^{l-2} a_{\alpha_k'}^{\dagger} a_{\alpha_l'}^{\dagger} a_{\alpha_l} a_{\alpha_k}
		|\alpha_k\alpha_l \underbrace{\alpha_1\alpha_2\dots\alpha_n}_{\neq \alpha_k,\alpha_l}\rangle \nonumber \\
	&=& (-1)^{k-1} (-1)^{l-2} 
	|\alpha_k'\alpha_l' \underbrace{\alpha_1\alpha_2\dots\alpha_n}_{\neq \alpha_k',\alpha_l'}\rangle \nonumber \\
	&=& |\alpha_1\alpha_2\dots\alpha_k'\dots\alpha_l'\dots\alpha_n\rangle \label{eq:2-35}
\end{align}

Inserting this in (\ref{eq:2-34}) gives
\begin{align}
	H_I |\alpha_1\alpha_2\dots\alpha_n\rangle
	&=& \sum_{\alpha_1', \alpha_2'} \langle \alpha_1'\alpha_2'|\hat{v}|\alpha_1\alpha_2\rangle
		a_{\alpha_1'}^{\dagger} a_{\alpha_2'}^{\dagger} a_{\alpha_2} a_{\alpha_1}
		|\alpha_1\alpha_2\dots\alpha_n\rangle \nonumber \\
	&+& \dots \nonumber \\
	&=& \sum_{\alpha_1', \alpha_n'} \langle \alpha_1'\alpha_n'|\hat{v}|\alpha_1\alpha_n\rangle
		a_{\alpha_1'}^{\dagger} a_{\alpha_n'}^{\dagger} a_{\alpha_n} a_{\alpha_1}
		|\alpha_1\alpha_2\dots\alpha_n\rangle \nonumber \\
	&+& \dots \nonumber \\
	&=& \sum_{\alpha_2', \alpha_n'} \langle \alpha_2'\alpha_n'|\hat{v}|\alpha_2\alpha_n\rangle
		a_{\alpha_2'}^{\dagger} a_{\alpha_n'}^{\dagger} a_{\alpha_n} a_{\alpha_2}
		|\alpha_1\alpha_2\dots\alpha_n\rangle \nonumber \\
	&+& \dots \nonumber \\
	&=& \sum_{\alpha, \beta, \gamma, \delta} ' \langle \alpha\beta|\hat{v}|\gamma\delta\rangle
		a^{\dagger}_\alpha a^{\dagger}_\beta a_\delta a_\gamma
		|\alpha_1\alpha_2\dots\alpha_n\rangle \label{eq:2-36}
\end{align}

Here we let $\sum'$ indicate that the sums running over $\alpha$ and $\beta$ run over all
single-particle states, while the summations  $\gamma$ and $\delta$ 
run over all pairs of single-particle states. We wish to remove this restriction and since
\begin{equation}
	\langle \alpha\beta|\hat{v}|\gamma\delta\rangle = \langle \beta\alpha|\hat{v}|\delta\gamma\rangle \label{eq:2-37}
\end{equation}
we get
\begin{align}
	\sum_{\alpha\beta} \langle \alpha\beta|\hat{v}|\gamma\delta\rangle a^{\dagger}_\alpha a^{\dagger}_\beta a_\delta a_\gamma &=& 
		\sum_{\alpha\beta} \langle \beta\alpha|\hat{v}|\delta\gamma\rangle 
		a^{\dagger}_\alpha a^{\dagger}_\beta a_\delta a_\gamma \label{eq:2-38a} \\
	&=& \sum_{\alpha\beta}\langle \beta\alpha|\hat{v}|\delta\gamma\rangle
		a^{\dagger}_\beta a^{\dagger}_\alpha a_\gamma a_\delta \label{eq:2-38b}
\end{align}
where we  have used the anti-commutation rules.

Changing the summation indices 
$\alpha$ and $\beta$ in (\ref{eq:2-38b}) we obtain
\begin{equation}
	\sum_{\alpha\beta} \langle \alpha\beta|\hat{v}|\gamma\delta\rangle a^{\dagger}_\alpha a^{\dagger}_\beta a_\delta a_\gamma =
		 \sum_{\alpha\beta} \langle \alpha\beta|\hat{v}|\delta\gamma\rangle 
		  a^{\dagger}_\alpha a^{\dagger}_\beta  a_\gamma a_\delta \label{eq:2-38c}
\end{equation}
From this it follows that the restriction on the summation over $\gamma$ and $\delta$ can be removed if we multiply with a factor $\frac{1}{2}$, resulting in 
\begin{equation}
	\hat{H}_I = \frac{1}{2} \sum_{\alpha\beta\gamma\delta} \langle \alpha\beta|\hat{v}|\gamma\delta\rangle
		a^{\dagger}_\alpha a^{\dagger}_\beta a_\delta a_\gamma \label{eq:2-39}
\end{equation}
where we sum freely over all single-particle states $\alpha$, 
$\beta$, $\gamma$ og $\delta$.

With this expression we can now verify that the second quantization form of $\hat{H}_I$ in Eq.~(\ref{eq:2-39}) 
results in the same matrix between two anti-symmetrized two-particle states as its corresponding coordinate
space representation. We have  
\begin{equation}
	\langle \alpha_1 \alpha_2|\hat{H}_I|\beta_1 \beta_2\rangle =
		\frac{1}{2} \sum_{\alpha\beta\gamma\delta}
			\langle \alpha\beta|\hat{v}|\gamma\delta\rangle \langle 0|a_{\alpha_2} a_{\alpha_1} 
			 a^{\dagger}_\alpha a^{\dagger}_\beta a_\delta a_\gamma 
			 a_{\beta_1}^{\dagger} a_{\beta_2}^{\dagger}|0\rangle. \label{eq:2-40}
\end{equation}

Using the commutation relations we get 
\begin{align}
	&& a_{\alpha_2} a_{\alpha_1}a^{\dagger}_\alpha a^{\dagger}_\beta 
		a_\delta a_\gamma a_{\beta_1}^{\dagger} a_{\beta_2}^{\dagger} \nonumber \\
	&=& a_{\alpha_2} a_{\alpha_1}a^{\dagger}_\alpha a^{\dagger}_\beta 
		( a_\delta \delta_{\gamma \beta_1} a_{\beta_2}^{\dagger} - 
		a_\delta  a_{\beta_1}^{\dagger} a_\gamma a_{\beta_2}^{\dagger} ) \nonumber \\
	&=& a_{\alpha_2} a_{\alpha_1}a^{\dagger}_\alpha a^{\dagger}_\beta 
		(\delta_{\gamma \beta_1} \delta_{\delta \beta_2} - \delta_{\gamma \beta_1} a_{\beta_2}^{\dagger} a_\delta -
		a_\delta a_{\beta_1}^{\dagger}\delta_{\gamma \beta_2} +
		a_\delta a_{\beta_1}^{\dagger} a_{\beta_2}^{\dagger} a_\gamma ) \nonumber \\
	&=& a_{\alpha_2} a_{\alpha_1}a^{\dagger}_\alpha a^{\dagger}_\beta 
		(\delta_{\gamma \beta_1} \delta_{\delta \beta_2} - \delta_{\gamma \beta_1} a_{\beta_2}^{\dagger} a_\delta \nonumber \\
		&& \qquad - \delta_{\delta \beta_1} \delta_{\gamma \beta_2} + \delta_{\gamma \beta_2} a_{\beta_1}^{\dagger} a_\delta
		+ a_\delta a_{\beta_1}^{\dagger} a_{\beta_2}^{\dagger} a_\gamma ) \label{eq:2-41}
\end{align}

The vacuum expectation value of this product of operators becomes
\begin{align}
	&& \langle 0|a_{\alpha_2} a_{\alpha_1} a^{\dagger}_\alpha a^{\dagger}_\beta a_\delta a_\gamma 
		a_{\beta_1}^{\dagger} a_{\beta_2}^{\dagger}|0\rangle \nonumber \\
	&=& (\delta_{\gamma \beta_1} \delta_{\delta \beta_2} -
		\delta_{\delta \beta_1} \delta_{\gamma \beta_2} ) 
		\langle 0|a_{\alpha_2} a_{\alpha_1}a^{\dagger}_\alpha a^{\dagger}_\beta|0\rangle \nonumber \\
	&=& (\delta_{\gamma \beta_1} \delta_{\delta \beta_2} -\delta_{\delta \beta_1} \delta_{\gamma \beta_2} )
	(\delta_{\alpha \alpha_1} \delta_{\beta \alpha_2} -\delta_{\beta \alpha_1} \delta_{\alpha \alpha_2} ) \label{eq:2-42b}
\end{align}

Insertion of 
Eq.~(\ref{eq:2-42b}) in Eq.~(\ref{eq:2-40}) results in
\begin{align}
	\langle \alpha_1\alpha_2|\hat{H}_I|\beta_1\beta_2\rangle &=& \frac{1}{2} \big[ 
		\langle \alpha_1\alpha_2|\hat{v}|\beta_1\beta_2\rangle- \langle \alpha_1\alpha_2|\hat{v}|\beta_2\beta_1\rangle \nonumber \\
		&& - \langle \alpha_2\alpha_1|\hat{v}|\beta_1\beta_2\rangle + \langle \alpha_2\alpha_1|\hat{v}|\beta_2\beta_1\rangle \big] \nonumber \\
	&=& \langle \alpha_1\alpha_2|\hat{v}|\beta_1\beta_2\rangle - \langle \alpha_1\alpha_2|\hat{v}|\beta_2\beta_1\rangle \nonumber \\
	&=& \langle \alpha_1\alpha_2|\hat{v}|\beta_1\beta_2\rangle_{\mathrm{AS}}. \label{eq:2-43b}
\end{align}

The two-body operator can also be expressed in terms of the anti-symmetrized matrix elements we discussed previously as
\begin{align}
	\hat{H}_I &=& \frac{1}{2} \sum_{\alpha\beta\gamma\delta}  \langle \alpha \beta|\hat{v}|\gamma \delta\rangle
		a_\alpha^{\dagger} a_\beta^{\dagger} a_\delta a_\gamma \nonumber \\
	&=& \frac{1}{4} \sum_{\alpha\beta\gamma\delta} \left[ \langle \alpha \beta|\hat{v}|\gamma \delta\rangle -
		\langle \alpha \beta|\hat{v}|\delta\gamma \rangle \right] 
		a_\alpha^{\dagger} a_\beta^{\dagger} a_\delta a_\gamma \nonumber \\
	&=& \frac{1}{4} \sum_{\alpha\beta\gamma\delta} \langle \alpha \beta|\hat{v}|\gamma \delta\rangle_{\mathrm{AS}}
		a_\alpha^{\dagger} a_\beta^{\dagger} a_\delta a_\gamma \label{eq:2-45}
\end{align}

The factors in front of the operator, either  $\frac{1}{4}$ or 
$\frac{1}{2}$ tells whether we use antisymmetrized matrix elements or not. 

We can now express the Hamiltonian operator for a many-fermion system  in the occupation basis representation
as  
\begin{equation}
	H = \sum_{\alpha, \beta} \langle \alpha|\hat{t}+\hat{u}_{\mathrm{ext}}|\beta\rangle a_\alpha^{\dagger} a_\beta + \frac{1}{4} \sum_{\alpha\beta\gamma\delta}
		\langle \alpha \beta|\hat{v}|\gamma \delta\rangle a_\alpha^{\dagger} a_\beta^{\dagger} a_\delta a_\gamma. \label{eq:2-46b}
\end{equation}
This is the form we will use in the rest of these lectures, assuming that we work with anti-symmetrized two-body matrix elements.

\section{Particle-hole formalism}

The index conventions used here -- $i,j,k,l$ for hole states below the Fermi
level, $a,b,c,d$ for particle states above it, and $p,q,r,s$ for either -- were
introduced in Section~\ref{sec:particlehole}.  We now give them their proper
setting: the reference determinant becomes a new vacuum, and the operators are
redefined with respect to it.


Second quantization is a useful and elegant formalism  for constructing many-body  states and 
quantum mechanical operators. One can express and translate many physical processes
into simple pictures such as Feynman diagrams. Expecation values of many-body states are also easily calculated.
However, although the equations are seemingly easy to set up, from  a practical point of view, that is
the solution of Schroedinger's equation, there is no particular gain.
The many-body equation is equally hard to solve, irrespective of representation. 
The cliche that 
there is no free lunch brings us down to earth again.  
Note however that a transformation to a particular
basis, for cases where the interaction obeys specific symmetries, can ease the solution of Schroedinger's equation. 

But there is at least one important case where second quantization comes to our rescue.
It is namely easy to introduce another reference state than the pure vacuum $|0\rangle $, where all single-particle states are active.
With many particles present it is often useful to introduce another reference state  than the vacuum state$|0\rangle $. We will label this state $|c\rangle$ ($c$ for core) and as we will see it can reduce 
considerably the complexity and thereby the dimensionality of the many-body problem. It allows us to sum up to infinite order specific many-body correlations.  The particle-hole representation is one of these handy representations. 

In the original particle representation these states are products of the creation operators  $a_{\alpha_i}^\dagger$ acting on the true vacuum $|0\rangle $.
Following Eq.~(\ref{eq:2-2}) we have 
\begin{align}
 |\alpha_1\alpha_2\dots\alpha_{n-1}\alpha_n\rangle &=& a_{\alpha_1}^\dagger a_{\alpha_2}^\dagger \dots
					a_{\alpha_{n-1}}^\dagger a_{\alpha_n}^\dagger |0\rangle  \label{eq:2-47a} \\
	|\alpha_1\alpha_2\dots\alpha_{n-1}\alpha_n\alpha_{n+1}\rangle &=&
		a_{\alpha_1}^\dagger a_{\alpha_2}^\dagger \dots a_{\alpha_{n-1}}^\dagger a_{\alpha_n}^\dagger
		a_{\alpha_{n+1}}^\dagger |0\rangle  \label{eq:2-47b} \\
	|\alpha_1\alpha_2\dots\alpha_{n-1}\rangle &=& a_{\alpha_1}^\dagger a_{\alpha_2}^\dagger \dots
		a_{\alpha_{n-1}}^\dagger |0\rangle  \label{eq:2-47c}
\end{align}

If we use Eq.~(\ref{eq:2-47a}) as our new reference state, we can simplify considerably the representation of 
this state
\begin{equation}
	|c\rangle  \equiv |\alpha_1\alpha_2\dots\alpha_{n-1}\alpha_n\rangle =
		a_{\alpha_1}^\dagger a_{\alpha_2}^\dagger \dots a_{\alpha_{n-1}}^\dagger a_{\alpha_n}^\dagger |0\rangle  \label{eq:2-48a}
\end{equation}
The new reference states for the $n+1$ and $n-1$ states can then be written as
\begin{align}
	|\alpha_1\alpha_2\dots\alpha_{n-1}\alpha_n\alpha_{n+1}\rangle &=& (-1)^n a_{\alpha_{n+1}}^\dagger |c\rangle 
		\equiv (-1)^n |\alpha_{n+1}\rangle_c \label{eq:2-48b} \\
	|\alpha_1\alpha_2\dots\alpha_{n-1}\rangle &=& (-1)^{n-1} a_{\alpha_n} |c\rangle  
		\equiv (-1)^{n-1} |\alpha_{n-1}\rangle_c \label{eq:2-48c} 
\end{align}

The first state has one additional particle with respect to the new vacuum state
$|c\rangle $  and is normally referred to as a one-particle state or one particle added to the 
many-body reference state. 
The second state has one particle less than the reference vacuum state  $|c\rangle $ and is referred to as
a one-hole state. 
When dealing with a new reference state it is often convenient to introduce 
new creation and annihilation operators since we have 
from Eq.~(\ref{eq:2-48c})
\begin{equation}
	a_\alpha |c\rangle  \neq 0 \label{eq:2-49}
\end{equation}
since  $\alpha$ is contained  in $|c\rangle $, while for the true vacuum we have 
$a_\alpha |0\rangle  = 0$ for all $\alpha$.

The new reference state leads to the definition of new creation and annihilation operators
which satisfy the following relations
\begin{align}
	b_\alpha |c\rangle  &=& 0 \label{eq:2-50a} \\
	\{b_\alpha^\dagger , b_\beta^\dagger \} = \{b_\alpha , b_\beta \} &=& 0 \nonumber  \\
	\{b_\alpha^\dagger , b_\beta \} &=& \delta_{\alpha \beta} \label{eq:2-50c}
\end{align}
We assume also that the new reference state is properly normalized
\begin{equation}
	\langle c | c \rangle = 1 \label{eq:2-51}
\end{equation}

The physical interpretation of these new operators is that of so-called quasiparticle states.
This means that a state defined by the addition of one extra particle to a reference state $|c\rangle $ may not necesseraly be interpreted as one particle coupled to a core.
We define now new creation operators that act on a state $\alpha$ creating a new quasiparticle state
\begin{equation}
	b_\alpha^\dagger|c\rangle  = \Bigg\{ \begin{array}{ll}
		a_\alpha^\dagger |c\rangle  = |\alpha\rangle, & \alpha > F \\
		\\
		a_\alpha |c\rangle  = |\alpha^{-1}\rangle, & \alpha \leq F
	\end{array} \label{eq:2-52}
\end{equation}
where $F$ is the Fermi level representing the last  occupied single-particle orbit 
of the new reference state $|c\rangle $. 

The annihilation is the hermitian conjugate of the creation operator
\[
	b_\alpha = (b_\alpha^\dagger)^\dagger,
\]
resulting in
\begin{equation}
	b_\alpha^\dagger = \Bigg\{ \begin{array}{ll}
		a_\alpha^\dagger & \alpha > F \\
		\\
		a_\alpha & \alpha \leq F
	\end{array} \qquad 
	b_\alpha = \Bigg\{ \begin{array}{ll}
		a_\alpha & \alpha > F \\
		\\
		 a_\alpha^\dagger & \alpha \leq F
	\end{array} \label{eq:2-54}
\end{equation}

With the new creation and annihilation operator  we can now construct 
many-body quasiparticle states, with one-particle-one-hole states, two-particle-two-hole
states etc in the same fashion as we previously constructed many-particle states. 
We can write a general particle-hole state as
\begin{equation}
	|\beta_1\beta_2\dots \beta_{n_p} \gamma_1^{-1} \gamma_2^{-1} \dots \gamma_{n_h}^{-1}\rangle \equiv
		\underbrace{b_{\beta_1}^\dagger b_{\beta_2}^\dagger \dots b_{\beta_{n_p}}^\dagger}_{>F}
		\underbrace{b_{\gamma_1}^\dagger b_{\gamma_2}^\dagger \dots b_{\gamma_{n_h}}^\dagger}_{\leq F} |c\rangle \label{eq:2-56}
\end{equation}
We can now rewrite our one-body and two-body operators in terms of the new creation and annihilation operators.
The number operator becomes
\begin{equation}
	\hat{N} = \sum_\alpha a_\alpha^\dagger a_\alpha= 
\sum_{\alpha > F} b_\alpha^\dagger b_\alpha + n_c - \sum_{\alpha \leq F} b_\alpha^\dagger b_\alpha \label{eq:2-57b}
\end{equation}
where $n_c$ is the number of particle in the new vacuum state $|c\rangle $.  
The action of $\hat{N}$ on a many-body state results in 
\begin{equation}
	N |\beta_1\beta_2\dots \beta_{n_p} \gamma_1^{-1} \gamma_2^{-1} \dots \gamma_{n_h}^{-1}\rangle = (n_p + n_c - n_h) |\beta_1\beta_2\dots \beta_{n_p} \gamma_1^{-1} \gamma_2^{-1} \dots \gamma_{n_h}^{-1}\rangle \label{2-59}
\end{equation}
Here  $n=n_p +n_c - n_h$ is the total number of particles in the quasi-particle state of 
Eq.~(\ref{eq:2-56}). Note that  $\hat{N}$ counts the total number of particles  present 
\begin{equation}
	N_{qp} = \sum_\alpha b_\alpha^\dagger b_\alpha, \label{eq:2-60}
\end{equation}
gives us the number of quasi-particles as can be seen by computing
\begin{equation}
	N_{qp}= |\beta_1\beta_2\dots \beta_{n_p} \gamma_1^{-1} \gamma_2^{-1} \dots \gamma_{n_h}^{-1}\rangle
		= (n_p + n_h)|\beta_1\beta_2\dots \beta_{n_p} \gamma_1^{-1} \gamma_2^{-1} \dots \gamma_{n_h}^{-1}\rangle \label{eq:2-61}
\end{equation}
where $n_{qp} = n_p + n_h$ is the total number of quasi-particles.

We express the one-body operator $\hat{H}_0$ in terms of the quasi-particle creation and annihilation operators, resulting in
\begin{align}
	\hat{H}_0 &=& \sum_{\alpha\beta > F} \langle \alpha|\hat{h}_0|\beta\rangle  b_\alpha^\dagger b_\beta +
		\sum_{\alpha > F, \beta \leq F } \left[\langle \alpha|\hat{h}_0|\beta\rangle b_\alpha^\dagger b_\beta^\dagger + \langle \beta|\hat{h}_0|\alpha\rangle b_\beta  b_\alpha \right] \nonumber \\
	&+& \sum_{\alpha \leq F} \langle \alpha|\hat{h}_0|\alpha\rangle - \sum_{\alpha\beta \leq F} \langle \beta|\hat{h}_0|\alpha\rangle b_\alpha^\dagger b_\beta \label{eq:2-63b}
\end{align}
The first term  gives contribution only for particle states, while the last one
contributes only for holestates. The second term can create or destroy a set of
quasi-particles and 
the third term is the contribution  from the vacuum state $|c\rangle$.

Before we continue with the expressions for the two-body operator, we introduce a nomenclature we will use for the rest of this
text. It is inspired by the notation used in quantum chemistry.
We reserve the labels $i,j,k,\dots$ for hole states and $a,b,c,\dots$ for states above $F$, viz.~particle states.
This means also that we will skip the constraint $\leq F$ or $> F$ in the summation symbols. 
Our operator $\hat{H}_0$  reads now 
\begin{align}
	\hat{H}_0 &=& \sum_{ab} \langle a|\hat{h}|b\rangle b_a^\dagger b_b +
		\sum_{ai} \left[
		\langle a|\hat{h}|i\rangle b_a^\dagger b_i^\dagger + 
		\langle i|\hat{h}|a\rangle b_i  b_a \right] \nonumber \\
	&+& \sum_{i} \langle i|\hat{h}|i\rangle - 
		\sum_{ij} \langle j|\hat{h}|i\rangle
		b_i^\dagger b_j \label{eq:2-63c}
\end{align} 

The two-particle operator in the particle-hole formalism  is more complicated since we have
to translate four indices $\alpha\beta\gamma\delta$ to the possible combinations of particle and hole
states.  When performing the commutator algebra we can regroup the operator in five different terms
\begin{equation}
	\hat{H}_I = \hat{H}_I^{(a)} + \hat{H}_I^{(b)} + \hat{H}_I^{(c)} + \hat{H}_I^{(d)} + \hat{H}_I^{(e)} \label{eq:2-65}
\end{equation}
Using anti-symmetrized  matrix elements, 
bthe term  $\hat{H}_I^{(a)}$ is  
\begin{equation}
	\hat{H}_I^{(a)} = \frac{1}{4}
	\sum_{abcd} \langle ab|\hat{V}|cd\rangle 
		b_a^\dagger b_b^\dagger b_d b_c \label{eq:2-66}
\end{equation}

The next term $\hat{H}_I^{(b)}$  reads
\begin{equation}
	 \hat{H}_I^{(b)} = \frac{1}{4} \sum_{abci}\left(\langle ab|\hat{V}|ci\rangle b_a^\dagger b_b^\dagger b_i^\dagger b_c +\langle ai|\hat{V}|cb\rangle b_a^\dagger b_i b_b b_c\right) \label{eq:2-67b}
\end{equation}
This term conserves the number of quasiparticles but creates or removes a 
three-particle-one-hole  state. 
For $\hat{H}_I^{(c)}$  we have
\begin{align}
	\hat{H}_I^{(c)}& =& \frac{1}{4}
		\sum_{abij}\left(\langle ab|\hat{V}|ij\rangle b_a^\dagger b_b^\dagger b_j^\dagger b_i^\dagger +
		\langle ij|\hat{V}|ab\rangle b_a  b_b b_j b_i \right)+  \nonumber \\
	&&	\frac{1}{2}\sum_{abij}\langle ai|\hat{V}|bj\rangle b_a^\dagger b_j^\dagger b_b b_i + 
		\frac{1}{2}\sum_{abi}\langle ai|\hat{V}|bi\rangle b_a^\dagger b_b. \label{eq:2-68c}
\end{align}

The first line stands for the creation of a two-particle-two-hole state, while the second line represents
the creation to two one-particle-one-hole pairs
while the last term represents a contribution to the particle single-particle energy
from the hole states, that is an interaction between the particle states and the hole states
within the new vacuum  state.
The fourth term reads
\begin{align}
	 \hat{H}_I^{(d)}& = &\frac{1}{4} 
	 	\sum_{aijk}\left(\langle ai|\hat{V}|jk\rangle b_a^\dagger b_k^\dagger b_j^\dagger b_i+
\langle ji|\hat{V}|ak\rangle b_k^\dagger b_j b_i b_a\right)+\nonumber \\
&&\frac{1}{4}\sum_{aij}\left(\langle ai|\hat{V}|ji\rangle b_a^\dagger b_j^\dagger+
\langle ji|\hat{V}|ai\rangle - \langle ji|\hat{V}|ia\rangle b_j b_a \right). \label{eq:2-69d} 
\end{align}
The terms in the first line  stand for the creation of a particle-hole state 
interacting with hole states, we will label this as a two-hole-one-particle contribution. 
The remaining terms are a particle-hole state interacting with the holes in the vacuum state. 
Finally we have 
\begin{equation}
	\hat{H}_I^{(e)} = \frac{1}{4}
		 \sum_{ijkl}
		 \langle kl|\hat{V}|ij\rangle b_i^\dagger b_j^\dagger b_l b_k+
	        \frac{1}{2}\sum_{ijk}\langle ij|\hat{V}|kj\rangle b_k^\dagger b_i
	        +\frac{1}{2}\sum_{ij}\langle ij|\hat{V}|ij\rangle \label{eq:2-70d}
\end{equation}
The first terms represents the 
interaction between two holes while the second stands for the interaction between a hole and the remaining holes in the vacuum state.
It represents a contribution to single-hole energy  to first order.
The last term collects all contributions to the energy of the ground state of a closed-shell system arising
from hole-hole correlations.

\section{Summarizing and defining a normal-ordered Hamiltonian}

Everything collected here has an exact counterpart in
Chapter~\ref{chap:mbphys}.  The Slater determinant is
Eq.~(\ref{eq:SlatDet}); the antisymmetrised two-body matrix element is
Eq.~(\ref{eq:2-antisymmatrix}); and the energy functional that follows from
the normal-ordered Hamiltonian is Eq.~(\ref{eq:EHF}), whose minimisation gives
the Hartree-Fock equations~(\ref{eq:2-hfequations}).  The difference is that
the derivations there ran over permutations, while here they will run over
contractions.


\[
  \Phi_{AS}(\alpha_1, \dots, \alpha_A; x_1, \dots x_A)=
            \frac{1}{\sqrt{A}} \sum_{\hat{P}} (-1)^P \hat{P} \prod_{i=1}^A \psi_{\alpha_i}(x_i),
\]
which is equivalent with $|\alpha_1 \dots \alpha_A\rangle= a_{\alpha_1}^{\dagger} \dots a_{\alpha_A}^{\dagger} |0\rangle$. We have also
    \[
        a_p^\dagger|0\rangle = |p\rangle, \quad a_p |q\rangle = \delta_{pq}|0\rangle
    \]
\[
  \delta_{pq} = \left\{a_p, a_q^\dagger \right\},
\]
and 
\[
0 = \left\{a_p^\dagger, a_q \right\} = \left\{a_p, a_q \right\} = \left\{a_p^\dagger, a_q^\dagger \right\}
\]
\[
|\Phi_0\rangle = |\alpha_1 \dots \alpha_A\rangle, \quad \alpha_1, \dots, \alpha_A \leq \alpha_F
\]

\[
\left\{a_p^\dagger, a_q \right\}= \delta_{pq}, p, q \leq \alpha_F 
\]
\[
\left\{a_p, a_q^\dagger \right\} = \delta_{pq}, p, q > \alpha_F
\]
with         $i,j,\ldots \leq \alpha_F, \quad a,b,\ldots > \alpha_F, \quad p,q, \ldots - \textrm{any}$
\[
        a_i|\Phi_0\rangle = |\Phi_i\rangle, \hspace{0.5cm} a_a^\dagger|\Phi_0\rangle = |\Phi^a\rangle
\]
and         
\[
a_i^\dagger|\Phi_0\rangle = 0 \hspace{0.5cm}  a_a|\Phi_0\rangle = 0
\]

The one-body operator is defined as
\[
 \hat{F} = \sum_{pq} \langle p|\hat{f}|q\rangle a_p^\dagger a_q
\]
while the two-body opreator is defined as
\[
\hat{V} = \frac{1}{4} \sum_{pqrs} \langle pq|\hat{v}|rs\rangle_{AS} a_p^\dagger a_q^\dagger a_s a_r
\]
where we have defined the antisymmetric matrix elements
\[
\langle pq|\hat{v}|rs\rangle_{AS} = \langle pq|\hat{v}|rs\rangle - \langle pq|\hat{v}|sr\rangle.
\]

We can also define a three-body operator
\[
\hat{V}_3 = \frac{1}{36} \sum_{pqrstu} \langle pqr|\hat{v}_3|stu\rangle_{AS} 
                a_p^\dagger a_q^\dagger a_r^\dagger a_u a_t a_s
\]
with the antisymmetrized matrix element
\begin{align}
            \langle pqr|\hat{v}_3|stu\rangle_{AS} = \langle pqr|\hat{v}_3|stu\rangle + \langle pqr|\hat{v}_3|tus\rangle + \langle pqr|\hat{v}_3|ust\rangle- \langle pqr|\hat{v}_3|sut\rangle - \langle pqr|\hat{v}_3|tsu\rangle - \langle pqr|\hat{v}_3|uts\rangle.
\end{align}


%----------------------------------------------------------------------
\section{Normal ordering}
\label{sec:normalordering}

\index{normal ordering}
The manipulations of the previous sections all follow the same pattern.  Take
a product of creation and annihilation operators, push the annihilation
operators to the right using the anticommutation relations until they meet the
vacuum and give zero, and collect the Kronecker deltas that fall out along the
way.  For two operators this is immediate.  For the four operators of an
overlap between two-particle states it is already an exercise in care, and for
the eight or more operators that arise in perturbation theory and
coupled-cluster theory it is hopeless.  Wick's theorem organises the whole
calculation once and for all.  We need two definitions first.

A product of creation and annihilation operators is in \emph{normal-ordered}
form when all creation operators stand to the left of all annihilation
operators.  For an arbitrary product we define the normal-ordered product
\begin{equation}
  N\left[xyz\cdots w\right] = (-1)^{\kappa}\,
    \Big(\text{the same operators, all creation operators to the left,}
    \Big)
  \label{eq:3-normalorder}
\end{equation}
where $\kappa$ counts the interchanges of neighbouring operators needed to
bring the product into that form.  Every interchange of two fermion operators
costs a minus sign, and the factor $(-1)^{\kappa}$ records them.  Two
elementary cases are
\begin{equation}
  N\left[a_p a_q^{\dagger}\right] = -a_q^{\dagger}a_p ,
  \qquad
  N\left[a_p^{\dagger} a_q\right] = a_p^{\dagger}a_q ,
  \label{eq:3-normalexamples}
\end{equation}
the first requiring one interchange and the second none.  The ordering of the
creation operators among themselves, and of the annihilation operators among
themselves, is a matter of convention as long as the same sign rule is applied
consistently.

The one property that makes normal ordering useful follows immediately from
$a_p\vert 0\rangle = 0$ and $\langle 0\vert a_p^{\dagger}=0$:
\begin{equation}
  \boxed{\;\langle 0\vert N\left[xyz\cdots w\right]\vert 0\rangle = 0\;}
  \qquad\text{for every non-empty product.}
  \label{eq:3-normalvanishes}
\end{equation}
Whatever else it contains, a normal-ordered product has an annihilation
operator on the far right or a creation operator on the far left, and either
one kills the vacuum.  A normal-ordered product therefore contributes nothing
to a vacuum expectation value.  The whole problem of evaluating
$\langle 0\vert xyz\cdots w\vert 0\rangle$ reduces to finding the pieces that
are \emph{not} normal-ordered.

%----------------------------------------------------------------------
\section{Contractions}
\label{sec:contractions}

\index{contraction}
The \emph{contraction} of two operators is defined as the difference between
their product and their normal-ordered product,
\begin{equation}
  \wstrut
  \wmark{ca}{x}\,\wmark{cb}{y}\wline[1.6ex]{ca}{cb}
   \;\equiv\; xy - N\left[xy\right] ,
  \label{eq:3-contractiondef}
\end{equation}
and we indicate it by a line joining the two operators.  Since the
normal-ordered piece has vanishing vacuum expectation value by
Eq.~(\ref{eq:3-normalvanishes}), a contraction is just a number,
\begin{equation}
  \wstrut
  \wmark{cc}{x}\,\wmark{cd}{y}\wline[1.6ex]{cc}{cd}
   = \langle 0\vert xy\vert 0\rangle .
  \label{eq:3-contractionvev}
\end{equation}
There are four possible pairs of operators, and only one of them survives:
\begin{align}
  \wstrut
  \wmark{c1}{a_p}\,\wmark{c2}{a_q^{\dagger}}\wline[1.6ex]{c1}{c2}
   &= \langle 0\vert a_p a_q^{\dagger}\vert 0\rangle = \delta_{pq},
  \label{eq:3-contraction1}\\[2ex]
  \wmark{c3}{a_p}\,\wmark{c4}{a_q}\wline[1.6ex]{c3}{c4}
   &= \langle 0\vert a_p a_q\vert 0\rangle = 0,
  \label{eq:3-contraction2}\\[2ex]
  \wmark{c5}{a_q^{\dagger}}\,\wmark{c6}{a_p}\wline[1.6ex]{c5}{c6}
   &= \langle 0\vert a_q^{\dagger} a_p\vert 0\rangle = 0,
  \label{eq:3-contraction3}\\[2ex]
  \wmark{c7}{a_p^{\dagger}}\,\wmark{c8}{a_q^{\dagger}}\wline[1.6ex]{c7}{c8}
   &= \langle 0\vert a_p^{\dagger} a_q^{\dagger}\vert 0\rangle = 0 .
  \label{eq:3-contraction4}
\end{align}
Only an annihilation operator standing to the \emph{left} of a creation
operator gives a non-vanishing contraction.  This single fact carries all the
bookkeeping that follows: it is the statement that the vacuum contains no
particles, and every zero in the calculations below can be traced back to it.

When the two contracted operators are not adjacent, the contraction carries
the sign of the interchanges needed to bring them together.  For a set of
contractions the total sign is
\begin{equation}
  (-1)^{\nu},
  \qquad
  \nu = \text{number of crossings of the contraction lines},
  \label{eq:3-crossingsign}
\end{equation}
since each crossing corresponds to one interchange of neighbouring fermion
operators.  We shall use this rule constantly, and it is verified numerically
for arbitrary strings in \texttt{wick.py}.

%----------------------------------------------------------------------
\section{Two operators}
\label{sec:wicktwo}

We now build the theorem up from the simplest cases, using nothing more than
the anticommutation relations.  For two operators the content of the theorem
is the definition~(\ref{eq:3-contractiondef}) rearranged,
\begin{equation}
  \wstrut
  xy = N\left[xy\right] +
  \wmark{w2a}{x}\,\wmark{w2b}{y}\wline[1.6ex]{w2a}{w2b} .
  \label{eq:3-wick2}
\end{equation}
Take $x=a_p$ and $y=a_q^{\dagger}$.  The anticommutation relation
$\{a_p,a_q^{\dagger}\}=\delta_{pq}$ gives
\begin{equation}
  a_p a_q^{\dagger} = \delta_{pq} - a_q^{\dagger}a_p
                    = \delta_{pq} + N\left[a_p a_q^{\dagger}\right],
  \label{eq:3-wick2example}
\end{equation}
where the second step is Eq.~(\ref{eq:3-normalexamples}).  Comparing with
Eq.~(\ref{eq:3-wick2}) identifies the contraction as $\delta_{pq}$, which is
Eq.~(\ref{eq:3-contraction1}).  Taking the vacuum expectation value and using
Eq.~(\ref{eq:3-normalvanishes}) recovers the overlap of two single-particle
states,
\begin{equation}
  \langle p\vert q\rangle
   = \langle 0\vert a_p a_q^{\dagger}\vert 0\rangle = \delta_{pq} .
  \label{eq:3-overlap1}
\end{equation}
The three remaining combinations $a_pa_q$, $a_p^{\dagger}a_q$ and
$a_p^{\dagger}a_q^{\dagger}$ are already normal-ordered up to a sign, so their
contractions vanish and so do their vacuum expectation values.

%----------------------------------------------------------------------
\section{Three operators}
\label{sec:wickthree}

An odd number of operators can never be paired up completely.  Whatever
contractions we carry out, at least one operator is always left over, and a
single uncontracted operator inside a vacuum expectation value gives
\begin{equation}
  \langle 0\vert a_p\vert 0\rangle
   = \langle 0\vert a_p^{\dagger}\vert 0\rangle = 0 .
  \label{eq:3-oddvanishes}
\end{equation}
Hence
\begin{equation}
  \langle 0\vert xyz\vert 0\rangle = 0
  \label{eq:3-wick3}
\end{equation}
for every choice of the three operators, as one may also check directly by
anticommuting.  Take $x=a_p$, $y=a_q^{\dagger}$, $z=a_r^{\dagger}$:
\begin{equation}
  a_p a_q^{\dagger}a_r^{\dagger}
   = \left(\delta_{pq}-a_q^{\dagger}a_p\right)a_r^{\dagger}
   = \delta_{pq}a_r^{\dagger}
   - a_q^{\dagger}\left(\delta_{pr}-a_r^{\dagger}a_p\right),
  \label{eq:3-three1}
\end{equation}
and every surviving term contains a single unpaired creation operator acting
on the vacuum, so the vacuum expectation value vanishes.

The result is worth stating explicitly because it is used constantly:
\emph{any} vacuum expectation value of an odd number of creation and
annihilation operators is zero.  The same argument reappears in the proof of
the theorem, where it disposes of half the terms in the induction step.

%----------------------------------------------------------------------
\section{Four operators}
\label{sec:wickfour}

Four operators is the first genuinely interesting case, and it is the one that
appears whenever we compute the overlap of two two-particle states,
\begin{equation}
  \vert rs\rangle = a_r^{\dagger}a_s^{\dagger}\vert 0\rangle,
  \qquad
  \vert pq\rangle = a_p^{\dagger}a_q^{\dagger}\vert 0\rangle,
  \qquad
  \langle rs\vert pq\rangle
   = \langle 0\vert a_s a_r a_p^{\dagger}a_q^{\dagger}\vert 0\rangle .
  \label{eq:3-fouroverlap}
\end{equation}

\paragraph{By anticommutation.}
We evaluate it the elementary way first, so that the answer can be checked
against the theorem afterwards.  Anticommuting $a_r$ past $a_p^{\dagger}$
with $a_r a_p^{\dagger}=\delta_{rp}-a_p^{\dagger}a_r$,
\begin{equation}
  a_s a_r a_p^{\dagger}a_q^{\dagger}
   = a_s\left(\delta_{rp}-a_p^{\dagger}a_r\right)a_q^{\dagger}
   = \delta_{rp}\,a_s a_q^{\dagger} - a_s a_p^{\dagger}a_r a_q^{\dagger} .
  \label{eq:3-four1}
\end{equation}
The first term is handled by Eq.~(\ref{eq:3-wick2example}),
\begin{equation}
  \delta_{rp}\,a_s a_q^{\dagger}
   = \delta_{rp}\left(\delta_{sq}-a_q^{\dagger}a_s\right)
   \;\longrightarrow\; \delta_{rp}\delta_{sq},
  \label{eq:3-four2}
\end{equation}
the arrow denoting that we take the vacuum expectation value and drop the
normal-ordered remainder.  For the second term we anticommute twice more,
\begin{equation}
  -a_s a_p^{\dagger}a_r a_q^{\dagger}
   = -a_s a_p^{\dagger}\left(\delta_{rq}-a_q^{\dagger}a_r\right)
   = -\delta_{rq}\,a_s a_p^{\dagger} + a_s a_p^{\dagger}a_q^{\dagger}a_r ,
  \label{eq:3-four3}
\end{equation}
and the last term has an annihilation operator on the far right, so it dies on
the vacuum.  The surviving piece is
\begin{equation}
  -\delta_{rq}\,a_s a_p^{\dagger}
   = -\delta_{rq}\left(\delta_{sp}-a_p^{\dagger}a_s\right)
   \;\longrightarrow\; -\delta_{rq}\delta_{sp} .
  \label{eq:3-four4}
\end{equation}
Collecting the two contributions,
\begin{equation}
  \boxed{\;
  \langle rs\vert pq\rangle
   = \langle 0\vert a_s a_r a_p^{\dagger}a_q^{\dagger}\vert 0\rangle
   = \delta_{rp}\delta_{sq} - \delta_{sp}\delta_{rq} \;}
  \label{eq:3-fourresult}
\end{equation}
Two-particle states are orthonormal up to precisely the antisymmetry expected
of fermions: the overlap is the antisymmetrised product of single-particle
overlaps.

\paragraph{By contractions.}
The same answer can be read off without any algebra at all.  Four operators
can be paired in three ways.  The first pairs the two annihilation operators
with each other and the two creation operators with each other,
\begin{equation}
  \wstrut
  \wmark{p1a}{a_s}\,\wmark{p1b}{a_r}\,
  \wmark{p1c}{a_p^{\dagger}}\,\wmark{p1d}{a_q^{\dagger}}
  \wline[1.6ex]{p1a}{p1b}\wline[1.6ex]{p1c}{p1d}
  \;=\;0 ,
  \label{eq:3-pair1}
\end{equation}
which vanishes by Eqs.~(\ref{eq:3-contraction2}) and
(\ref{eq:3-contraction4}).  The second is the nested pairing, with no
crossings and hence a plus sign,
\begin{equation}
  \wstrut
  \wmark{p2a}{a_s}\,\wmark{p2b}{a_r}\,
  \wmark{p2c}{a_p^{\dagger}}\,\wmark{p2d}{a_q^{\dagger}}
  \wline[3.4ex]{p2a}{p2d}\wline[1.4ex]{p2b}{p2c}
  \;=\;+\,\delta_{sq}\delta_{rp} .
  \label{eq:3-pair2}
\end{equation}
The third has the two lines crossing once, and therefore a minus sign by
Eq.~(\ref{eq:3-crossingsign}); we draw one line above and one below,
\begin{equation}
  \wstrut
  \wmark{p3a}{a_s}\,\wmark{p3b}{a_r}\,
  \wmark{p3c}{a_p^{\dagger}}\,\wmark{p3d}{a_q^{\dagger}}
  \wline[1.6ex]{p3a}{p3c}\wlineb[1.6ex]{p3b}{p3d}
  \;=\;-\,\delta_{sp}\delta_{rq} .
  \label{eq:3-pair3}
\end{equation}
Adding the three gives Eq.~(\ref{eq:3-fourresult}) again.

Three pairings, one of which vanishes on inspection and two of which are read
straight off: this is the economy the theorem buys.  The program
\texttt{wick.py} carries out both routes for arbitrary strings and confirms
that they agree.  For the eight-operator strings of
Section~\ref{sec:wickapps} the anticommutation recursion takes $129$ steps,
while the theorem enumerates $105$ pairings of which only a handful survive.

%----------------------------------------------------------------------
\section{Wick's theorem}
\label{sec:wicktheorem}

\index{Wick's theorem}
The pattern of the previous three sections generalises.  \emph{Wick's
theorem} states that an arbitrary product of creation and annihilation
operators can be written as the normal-ordered product plus the sum of all
possible contractions,
\begin{align}
  xyz\cdots w
   &= N\left[xyz\cdots w\right] \nonumber\\
   &\quad + \sum_{(1)} N\left[xyz\cdots w\right]
     \qquad\text{(all single contractions)}\nonumber\\
   &\quad + \sum_{(2)} N\left[xyz\cdots w\right]
     \qquad\text{(all double contractions)}\nonumber\\
   &\quad + \cdots
     + \sum_{[M/2]} N\left[xyz\cdots w\right] ,
  \label{eq:3-wicktheorem}
\end{align}
where $\sum_{(m)}$ runs over all ways of choosing $m$ disjoint pairs of
operators to contract, each term carrying the sign~(\ref{eq:3-crossingsign}),
and $[M/2]$ is the largest number of disjoint pairs that $M$ operators admit.

The theorem becomes a computational tool the moment we take a vacuum
expectation value.  Every term in Eq.~(\ref{eq:3-wicktheorem}) that still
contains an uncontracted operator is normal-ordered and therefore vanishes by
Eq.~(\ref{eq:3-normalvanishes}).  Only the \emph{fully} contracted terms
survive:
\begin{equation}
  \boxed{\;
  \langle 0\vert xyz\cdots w\vert 0\rangle
   = \sum_{[M/2]} N\left[xyz\cdots w\right]
   \quad\text{(fully contracted terms only)} \;}
  \label{eq:3-wickvev}
\end{equation}
For $M$ even there are $(M-1)!!$ such terms -- one, three, fifteen and
one hundred and five for $M=2,4,6,8$ -- and most of them vanish on inspection
because one of their contractions is of the type
(\ref{eq:3-contraction2})--(\ref{eq:3-contraction4}).  For $M$ odd there are
none, which is Eq.~(\ref{eq:3-wick3}) again.

The practical consequence is the one anticipated in
Section~\ref{sec:expdim}.  Because an $n$-body operator can connect only
Slater determinants differing by at most $n$ single-particle states, a
configuration-interaction Hamiltonian built from one- and two-body operators
is extremely sparse: of the $\binom{n}{N}^2$ possible matrix elements only a
vanishing fraction are non-zero.  That sparsity is precisely what the Lanczos
algorithm of Section~\ref{sec:lanczos} exploits, and it is why the Hamiltonian
never has to be stored.

%----------------------------------------------------------------------
\section{Proof of Wick's theorem}
\label{sec:wickproof}

The proof is by induction on the number of operators.  The case $M=2$ is
Eq.~(\ref{eq:3-wick2}), established directly from the anticommutation
relation.  The induction step requires one preparatory result, which does all
the real work.

\paragraph{A lemma: appending one operator on the right.}
Let $x,y,z,\dots,w$ be a set of operators and let $\Omega$ be one further
creation or annihilation operator.  Then
\begin{equation}
  N\left[xyz\cdots w\right]\,\Omega
   = N\left[xyz\cdots w\,\Omega\right]
   + \sum_{(1)\Omega} N\left[xyz\cdots w\,\Omega\right] ,
  \label{eq:3-lemma}
\end{equation}
where $\sum_{(1)\Omega}$ runs over all single contractions \emph{involving}
$\Omega$, that is over the contractions of $\Omega$ with each of
$x,y,z,\dots,w$ in turn.

\paragraph{Proof of the lemma.}
There are four cases to consider, and three of them are immediate.

\begin{enumerate}
\item[(i)] \emph{$\Omega$ is an annihilation operator.}  Then
  $N[xyz\cdots w]\,\Omega$ is already in normal-ordered form, since appending
  an annihilation operator on the far right cannot disturb the ordering, and
  so the left-hand side equals $N[xyz\cdots w\,\Omega]$.  At the same time
  every contraction with $\Omega$ vanishes: a non-zero contraction requires
  an annihilation operator to the \emph{left} of a creation operator, and here
  $\Omega$ is itself an annihilation operator standing to the right of
  everything.  Both extra terms are absent and the lemma holds trivially.

\item[(ii)] \emph{All of $x,y,\dots,w$ are annihilation operators and $\Omega$
  is an annihilation operator.}  This is a special case of (i), and the lemma
  is again trivially satisfied.

\item[(iii)] \emph{$\Omega$ is a creation operator.}  Here we need only prove
  the lemma when all of $x,y,\dots,w$ are annihilation operators.  The reason
  is that a creation operator inside the normal-ordered product already
  stands to the left of $\Omega$'s eventual position, and its contraction with
  $\Omega$ vanishes by Eq.~(\ref{eq:3-contraction4}),
  \[
    \wmark{L1}{a_i^{\dagger}}\,\wmark{L2}{a_j^{\dagger}}\wline[1.6ex]{L1}{L2}=0 .
  \]
  Creation operators in the string thus contribute neither to the reordering
  nor to the sum of contractions, and may be carried along untouched.

\item[(iv)] It remains to anticommute $\Omega$ leftwards through the
  annihilation operators $x,y,\dots,w$, one at a time.  Each step uses
  \[
    a_j\,\Omega = \delta_{j\Omega} - \Omega\,a_j ,
  \]
  and produces two pieces: a Kronecker delta, which is precisely the
  contraction of $a_j$ with $\Omega$ carrying the sign of the operators
  $\Omega$ has had to cross, and a term in which $\Omega$ has moved one place
  to the left.  Iterating until $\Omega$ has passed all of them, the
  accumulated delta terms are exactly the sum
  $\sum_{(1)\Omega} N[xyz\cdots w\,\Omega]$, and the final term, in which
  $\Omega$ stands to the left of every annihilation operator, is exactly
  $N[xyz\cdots w\,\Omega]$.  This is Eq.~(\ref{eq:3-lemma}). \hfill$\square$
\end{enumerate}

\paragraph{The lemma at work.}
For $M=1$ the lemma reads
\begin{equation}
  \wstrut
  N\left[a_1\right]a_{\ell}^{\dagger}
   = N\left[a_1 a_{\ell}^{\dagger}\right]
   + N\left[\wmark{L3}{a_1}\,\wmark{L4}{a_{\ell}^{\dagger}}
       \wline[1.6ex]{L3}{L4}\right]
   = -\langle 0\vert a_{\ell}^{\dagger}a_1\vert 0\rangle
     + \delta_{1\ell} ,
  \label{eq:3-lemmaM1}
\end{equation}
which is just Eq.~(\ref{eq:3-wick2example}) once more.  For $M=2$,
\begin{equation}
  \wstrut
  a_0\,N\left[a_1\right]a_{\ell}^{\dagger}
   = N\left[a_0 a_1 a_{\ell}^{\dagger}\right]
   + N\left[a_0\,\wmark{L5}{a_1}\,\wmark{L6}{a_{\ell}^{\dagger}}
       \wline[1.6ex]{L5}{L6}\right]
   + N\left[\wmark{L7}{a_0}\,a_1\,\wmark{L8}{a_{\ell}^{\dagger}}
       \wline[2.8ex]{L7}{L8}\right],
  \label{eq:3-lemmaM2}
\end{equation}
that is $N[a_0a_1a_{\ell}^{\dagger}] + \delta_{1\ell}N[a_0]
+ \delta_{0\ell}N[a_1]$ up to the signs of the crossings.  Generally,
\begin{equation}
  a_0\,N\left[a_1a_2\cdots a_M\right]a_{\ell}^{\dagger}
   = \sum_{(1)\Omega} N\left[a_0a_1\cdots a_M a_{\ell}^{\dagger}\right]
   + N\left[a_0a_1\cdots a_M a_{\ell}^{\dagger}\right] .
  \label{eq:3-lemmaMgeneral}
\end{equation}

\paragraph{The induction step.}
Assume Wick's theorem, Eq.~(\ref{eq:3-wicktheorem}), holds for a product of
$M$ operators.  Multiply both sides from the right by one further operator
$\Omega$ and apply the lemma to every term.  The first term gives
\begin{equation}
  N\left[xyz\cdots w\right]\Omega
   = N\left[xyz\cdots w\,\Omega\right]
   + \sum_{(1)\Omega} N\left[xyz\cdots w\,\Omega\right],
  \label{eq:3-induct1}
\end{equation}
the $m$-fold contracted terms give
\begin{equation}
  \sum_{(m)} N\left[xyz\cdots w\right]\Omega
   = \sum_{(m)\ne\Omega} N\left[xyz\cdots w\,\Omega\right]
   + \sum_{(m)\Omega} N\left[xyz\cdots w\,\Omega\right],
  \label{eq:3-induct2}
\end{equation}
where $\sum_{(m)\ne\Omega}$ collects the $m$-fold contractions that do not
involve $\Omega$ and $\sum_{(m)\Omega}$ those in which one of the $m$
contractions ends on $\Omega$.  But an $m$-fold contraction of the $M+1$
operators either involves $\Omega$ or it does not, so the two sums recombine,
\begin{equation}
  \sum_{(m)\ne\Omega} + \sum_{(m)\Omega} = \sum_{(m)}
  \qquad\text{over all } M+1 \text{ operators.}
  \label{eq:3-recombine}
\end{equation}
Collecting the terms by the number of contractions therefore reproduces
Eq.~(\ref{eq:3-wicktheorem}) with $M$ replaced by $M+1$,
\begin{equation}
  xyz\cdots w\,\Omega
   = N\left[xyz\cdots w\,\Omega\right]
   + \sum_{(1)} N\left[\cdots\right]
   + \sum_{(2)} N\left[\cdots\right]
   + \cdots
   + \sum_{[(M+1)/2]} N\left[\cdots\right] .
  \label{eq:3-wickM1}
\end{equation}
This completes the induction. \hfill$\square$

\paragraph{A remark on parity.}
If $M+1$ is odd, the last sum $\sum_{[(M+1)/2]}$ still leaves one operator
uncontracted, since an odd number of operators cannot be paired up completely.
On taking the vacuum expectation value these terms vanish by
Eq.~(\ref{eq:3-oddvanishes}), and we recover the result of
Section~\ref{sec:wickthree}: vacuum expectation values of odd strings are
zero.  For $M+1$ even the last sum consists of the fully contracted terms,
and Eq.~(\ref{eq:3-wickvev}) follows.

%----------------------------------------------------------------------
\section{Wick's generalised theorem}
\label{sec:wickgeneral}

\index{Wick's theorem!generalised}
In practice we rarely meet a bare product of operators.  What we meet is a
product of groups, each of which is \emph{already} in normal-ordered form:
the bra, the operator, and the ket.  Wick's generalised theorem states that
for such a product
\begin{equation}
  N\left[A_1A_2\cdots\right]\,N\left[B_1B_2\cdots\right]\,
  N\left[C_1C_2\cdots\right]
   = N\left[A_1A_2\cdots B_1B_2\cdots C_1C_2\cdots\right]
   + \sum_{\mathrm{all}} N\left[A_1A_2\cdots C_1C_2\cdots\right],
  \label{eq:3-wickgeneral}
\end{equation}
where the sum runs over all contractions \emph{between operators belonging to
different normal-ordered groups}.  Contractions within a single group are
absent, precisely because that group is already normal-ordered and its
internal contractions vanish by
Eqs.~(\ref{eq:3-contraction2})--(\ref{eq:3-contraction4}).

This is a substantial saving.  In the examples of the next section the naive
count of pairings runs into the hundreds, while the number of contractions
between different groups is a handful.

%----------------------------------------------------------------------
\section{Applications: matrix elements between two-particle states}
\label{sec:wickapps}

We close with the two calculations that motivated the whole construction.
Both use the generalised theorem, and in both the bra, the operator and the
ket are separately normal-ordered.

\paragraph{A one-body operator.}
With
\begin{equation}
  \hat{O}^{(1)} = \sum_{pq}\langle p\vert\hat{o}^{(1)}\vert q\rangle\,
    a_p^{\dagger}a_q ,
  \label{eq:3-onebodyop}
\end{equation}
and the two-particle states $\vert ij\rangle=a_i^{\dagger}a_j^{\dagger}\vert
0\rangle$, $\langle ij\vert = \langle 0\vert a_j a_i$, we need
\begin{equation}
  \langle ij\vert \hat{O}^{(1)}\vert k\ell\rangle
   = \sum_{pq}\langle p\vert\hat{o}^{(1)}\vert q\rangle\,
     \langle 0\vert\,
     \underbrace{a_j a_i}_{\text{bra}}\;
     \underbrace{a_p^{\dagger}a_q}_{\text{operator}}\;
     \underbrace{a_k^{\dagger}a_{\ell}^{\dagger}}_{\text{ket}}\,
     \vert 0\rangle .
  \label{eq:3-onebodyme}
\end{equation}
Only contractions between different groups contribute, and each must join an
annihilation operator to a creation operator standing to its right.  One such
term is
\begin{equation}
  \wstrut
  \wmark{o1}{a_j}\,\wmark{o2}{a_i}\,\wmark{o3}{a_p^{\dagger}}\,
  \wmark{o4}{a_q}\,\wmark{o5}{a_k^{\dagger}}\,\wmark{o6}{a_{\ell}^{\dagger}}
  \wline[1.4ex]{o2}{o3}\wline[1.4ex]{o4}{o5}\wlineb{o1}{o6}
  \;=\;\delta_{ip}\,\delta_{qk}\,\delta_{j\ell},
  \label{eq:3-onebodycontraction}
\end{equation}
which contributes $\delta_{j\ell}\,\langle i\vert\hat{o}^{(1)}\vert k\rangle$
after the sums over $p$ and $q$ are carried out.  Collecting the four such
terms gives the familiar result
\begin{equation}
  \langle ij\vert \hat{O}^{(1)}\vert k\ell\rangle
   = \delta_{j\ell}\langle i\vert\hat{o}^{(1)}\vert k\rangle
   - \delta_{jk}\langle i\vert\hat{o}^{(1)}\vert\ell\rangle
   - \delta_{i\ell}\langle j\vert\hat{o}^{(1)}\vert k\rangle
   + \delta_{ik}\langle j\vert\hat{o}^{(1)}\vert\ell\rangle .
  \label{eq:3-onebodyresult}
\end{equation}
The matrix element vanishes whenever the two Slater determinants differ by
more than one single-particle state, since then no assignment of the deltas
can be satisfied.  This is the Condon-Slater rule of
Section~\ref{sec:particlehole}, now obtained without a single explicit
permutation.

\paragraph{The two-body interaction.}
With the normal-ordered interaction
\begin{equation}
  \hat{H}_I = \frac{1}{2}\sum_{pqrs}\langle pq\vert v\vert rs\rangle\,
    a_p^{\dagger}a_q^{\dagger}a_s a_r ,
  \label{eq:3-twobodyop}
\end{equation}
the diagonal matrix element in the state $\vert ij\rangle$ is
\begin{equation}
  \langle ij\vert \hat{H}_I\vert ij\rangle
   = \frac{1}{2}\sum_{pqrs}\langle pq\vert v\vert rs\rangle\,
     \langle 0\vert\,
     \underbrace{a_j a_i}_{\text{bra}}\;
     \underbrace{a_p^{\dagger}a_q^{\dagger}a_s a_r}_{\text{operator}}\;
     \underbrace{a_i^{\dagger}a_j^{\dagger}}_{\text{ket}}\,
     \vert 0\rangle .
  \label{eq:3-twobodyme}
\end{equation}
This is an eight-operator string.  Enumerating its $105$ pairings by hand
would be unpleasant, but the generalised theorem restricts us to contractions
between the three groups, and of those only four survive.  Two of them are
\begin{equation}
  \wstrut
  \wmark{t1}{a_j}\,\wmark{t2}{a_i}\,\wmark{t3}{a_p^{\dagger}}\,
  \wmark{t4}{a_q^{\dagger}}\,\wmark{t5}{a_s}\,\wmark{t6}{a_r}\,
  \wmark{t7}{a_i^{\dagger}}\,\wmark{t8}{a_j^{\dagger}}
  \wline[1.4ex]{t2}{t3}\wline[1.4ex]{t6}{t7}
  \wlineb[1.4ex]{t1}{t4}\wlineb[3.0ex]{t5}{t8}
  \;=\;\delta_{ip}\delta_{jq}\delta_{ri}\delta_{sj},
  \label{eq:3-twobodyc1}
\end{equation}
giving $\langle ij\vert v\vert ij\rangle$, and the one obtained by exchanging
the roles of $r$ and $s$, which crosses one more line and therefore enters
with the opposite sign, giving $-\langle ij\vert v\vert ji\rangle$.  The two
remaining terms are these two with $i$ and $j$ interchanged; they are equal to
the first two and cancel the factor $\tfrac12$.  The result is
\begin{equation}
  \boxed{\;
  \langle ij\vert \hat{H}_I\vert ij\rangle
   = \langle ij\vert v\vert ij\rangle - \langle ij\vert v\vert ji\rangle
   = \langle ij\vert v\vert ij\rangle_{\mathrm{AS}} \;}
  \label{eq:3-twobodyresult}
\end{equation}
which is exactly the antisymmetrised matrix element obtained in
Chapter~\ref{chap:mbphys} by writing out the antisymmetriser and sorting
through permutations.  The two derivations agree, as they must; the point is
that this one required no permutations at all, only the observation that a
contraction is non-zero when an annihilation operator stands to the left of a
creation operator.

The program \texttt{wick.py} performs every calculation of these sections in
both ways -- by brute-force anticommutation and by summing contractions with
their signs -- and verifies that they agree, for strings of two, four, six and
eight operators.  It also reproduces Eq.~(\ref{eq:3-twobodyresult}) by
enumerating which of the $(p,q,r,s)$ assignments survive.

\section{Operators in second quantization}
\label{sec:bitrep}

What follows is the bridge from the formalism to a running code.  The
Hamiltonian matrix built here is the sparse, never-stored matrix that the
Lanczos algorithm of Section~\ref{sec:lanczos} was designed for: only the
product $\bm{H}\bm{v}$ is needed, and the Condon-Slater rules of
Section~\ref{sec:wickapps} say which of its elements can be non-zero.  The bit
representation below is how that product is evaluated in practice.


In the build-up of a shell-model or FCI code that is meant to tackle large dimensionalities
is the action of the Hamiltonian $\hat{H}$ on a
Slater determinant represented in second quantization as
\[
 |\alpha_1\dots \alpha_n\rangle = a_{\alpha_1}^{\dagger} a_{\alpha_2}^{\dagger} \dots a_{\alpha_n}^{\dagger} |0\rangle.
\]
The time consuming part stems from the action of the Hamiltonian
on the above determinant,
\[
\left(\sum_{\alpha\beta} \langle \alpha|t+u|\beta\rangle a_\alpha^{\dagger} a_\beta + \frac{1}{4} \sum_{\alpha\beta\gamma\delta}
                \langle \alpha \beta|\hat{v}|\gamma \delta\rangle a_\alpha^{\dagger} a_\beta^{\dagger} a_\delta a_\gamma\right)a_{\alpha_1}^{\dagger} a_{\alpha_2}^{\dagger} \dots a_{\alpha_n}^{\dagger} |0\rangle.
\]
A practically useful way to implement this action is to encode a Slater determinant as a bit pattern.

Assume that we have at our disposal $n$ different single-particle orbits
$\alpha_0,\alpha_2,\dots,\alpha_{n-1}$ and that we can distribute  among these orbits $N\le n$ particles.

A Slater  determinant can then be coded as an integer of $n$ bits. As an example, if we have $n=16$ single-particle states
$\alpha_0,\alpha_1,\dots,\alpha_{15}$ and $N=4$ fermions occupying the states $\alpha_3$, $\alpha_6$, $\alpha_{10}$ and $\alpha_{13}$
we could write this Slater determinant as  
\[
\Phi_{\Lambda} = a_{\alpha_3}^{\dagger} a_{\alpha_6}^{\dagger} a_{\alpha_{10}}^{\dagger} a_{\alpha_{13}}^{\dagger} |0\rangle.
\]
The unoccupied single-particle states have bit value $0$ while the occupied ones are represented by bit state $1$. 
In the binary notation we would write this   16 bits long integer as
\[
\begin{array}{cccccccccccccccc}
{\alpha_0}&{\alpha_1}&{\alpha_2}&{\alpha_3}&{\alpha_4}&{\alpha_5}&{\alpha_6}&{\alpha_7} & {\alpha_8} &{\alpha_9} & {\alpha_{10}} &{\alpha_{11}} &{\alpha_{12}} &{\alpha_{13}} &{\alpha_{14}} & {\alpha_{15}} \\
{0} & {0} &{0} &{1} &{0} &{0} &{1} &{0} &{0} &{0} &{1} &{0} &{0} &{1} &{0} & {0} \\
\end{array}
\]
which translates into the decimal number
\[
2^3+2^6+2^{10}+2^{13}=9288.
\]
We can thus encode a Slater determinant as a bit pattern.

With $N$ particles that can be distributed over $n$ single-particle states, the total number of Slater determinats (and defining thereby the dimensionality of the system) is
\[
\mathrm{dim}(\mathcal{H}) = \left(\begin{array}{c} n \\N\end{array}\right).
\]
The total number of bit patterns is $2^n$. 

We assume again that we have at our disposal $n$ different single-particle orbits
$\alpha_0,\alpha_2,\dots,\alpha_{n-1}$ and that we can distribute  among these orbits $N\le n$ particles.
The ordering among these states is important as it defines the order of the creation operators.
We will write the determinant 
\[
\Phi_{\Lambda} = a_{\alpha_3}^{\dagger} a_{\alpha_6}^{\dagger} a_{\alpha_{10}}^{\dagger} a_{\alpha_{13}}^{\dagger} |0\rangle,
\]
in a more compact way as 
\[
\Phi_{3,6,10,13} = |0001001000100100\rangle.
\]
The action of a creation operator is thus
\[
a^{\dagger}_{\alpha_4}\Phi_{3,6,10,13} = a^{\dagger}_{\alpha_4}|0001001000100100\rangle=a^{\dagger}_{\alpha_4}a_{\alpha_3}^{\dagger} a_{\alpha_6}^{\dagger} a_{\alpha_{10}}^{\dagger} a_{\alpha_{13}}^{\dagger} |0\rangle,
\]
which becomes
\[
-a_{\alpha_3}^{\dagger} a^{\dagger}_{\alpha_4} a_{\alpha_6}^{\dagger} a_{\alpha_{10}}^{\dagger} a_{\alpha_{13}}^{\dagger} |0\rangle=-|0001101000100100\rangle.
\]

Similarly
\[
a^{\dagger}_{\alpha_6}\Phi_{3,6,10,13} = a^{\dagger}_{\alpha_6}|0001001000100100\rangle=a^{\dagger}_{\alpha_6}a_{\alpha_3}^{\dagger} a_{\alpha_6}^{\dagger} a_{\alpha_{10}}^{\dagger} a_{\alpha_{13}}^{\dagger} |0\rangle,
\]
which becomes
\[
-a^{\dagger}_{\alpha_4} (a_{\alpha_6}^{\dagger})^ 2 a_{\alpha_{10}}^{\dagger} a_{\alpha_{13}}^{\dagger} |0\rangle=0!
\]
This gives a simple recipe:  
\begin{itemize}
\item If one of the bits $b_j$ is $1$ and we act with a creation operator on this bit, we return a null vector

\item If $b_j=0$, we set it to $1$ and return a sign factor $(-1)^l$, where $l$ is the number of bits set before bit $j$.
\end{itemize}

\noindent
Consider the action of $a^{\dagger}_{\alpha_2}$ on various slater determinants:
\[
\begin{array}{ccc}
a^{\dagger}_{\alpha_2}\Phi_{00111}& = a^{\dagger}_{\alpha_2}|00111\rangle&=0\times |00111\rangle\\
a^{\dagger}_{\alpha_2}\Phi_{01011}& = a^{\dagger}_{\alpha_2}|01011\rangle&=(-1)\times |01111\rangle\\
a^{\dagger}_{\alpha_2}\Phi_{01101}& = a^{\dagger}_{\alpha_2}|01101\rangle&=0\times |01101\rangle\\
a^{\dagger}_{\alpha_2}\Phi_{01110}& = a^{\dagger}_{\alpha_2}|01110\rangle&=0\times |01110\rangle\\
a^{\dagger}_{\alpha_2}\Phi_{10011}& = a^{\dagger}_{\alpha_2}|10011\rangle&=(-1)\times |10111\rangle\\
a^{\dagger}_{\alpha_2}\Phi_{10101}& = a^{\dagger}_{\alpha_2}|10101\rangle&=0\times |10101\rangle\\
a^{\dagger}_{\alpha_2}\Phi_{10110}& = a^{\dagger}_{\alpha_2}|10110\rangle&=0\times |10110\rangle\\
a^{\dagger}_{\alpha_2}\Phi_{11001}& = a^{\dagger}_{\alpha_2}|11001\rangle&=(+1)\times |11101\rangle\\
a^{\dagger}_{\alpha_2}\Phi_{11010}& = a^{\dagger}_{\alpha_2}|11010\rangle&=(+1)\times |11110\rangle\\
\end{array}
\]
What is the simplest way to obtain the phase when we act with one annihilation(creation) operator
on the given Slater determinant representation?

We have an SD representation
\[
\Phi_{\Lambda} = a_{\alpha_0}^{\dagger} a_{\alpha_3}^{\dagger} a_{\alpha_6}^{\dagger} a_{\alpha_{10}}^{\dagger} a_{\alpha_{13}}^{\dagger} |0\rangle,
\]
in a more compact way as
\[
\Phi_{0,3,6,10,13} = |1001001000100100\rangle.
\]
The action of
\[
a^{\dagger}_{\alpha_4}a_{\alpha_0}\Phi_{0,3,6,10,13} = a^{\dagger}_{\alpha_4}|0001001000100100\rangle=a^{\dagger}_{\alpha_4}a_{\alpha_3}^{\dagger} a_{\alpha_6}^{\dagger} a_{\alpha_{10}}^{\dagger} a_{\alpha_{13}}^{\dagger} |0\rangle,
\]
which becomes
\[
-a_{\alpha_3}^{\dagger} a^{\dagger}_{\alpha_4} a_{\alpha_6}^{\dagger} a_{\alpha_{10}}^{\dagger} a_{\alpha_{13}}^{\dagger} |0\rangle=-|0001101000100100\rangle.
\]

The action
\[
a_{\alpha_0}\Phi_{0,3,6,10,13} = |0001001000100100\rangle,
\]
can be obtained by subtracting the logical sum (AND operation) of $\Phi_{0,3,6,10,13}$ and 
a word which represents only $\alpha_0$, that is
\[
|1000000000000000\rangle,
\] 
from $\Phi_{0,3,6,10,13}= |1001001000100100\rangle$.

This operation gives $|0001001000100100\rangle$. 

Similarly, we can form $a^{\dagger}_{\alpha_4}a_{\alpha_0}\Phi_{0,3,6,10,13}$, say, by adding 
$|0000100000000000\rangle$ to $a_{\alpha_0}\Phi_{0,3,6,10,13}$, first checking that their logical sum
is zero in order to make sure that orbital $\alpha_4$ is not already occupied. 

It is trickier however to get the phase $(-1)^l$. 
One possibility is as follows
\begin{itemize}
\item Let $S_1$ be a word that represents the $1-$bit to be removed and all others set to zero.
\end{itemize}

\noindent
In the previous example $S_1=|1000000000000000\rangle$
\begin{itemize}
\item Define $S_2$ as the similar word that represents the bit to be added, that is in our case
\end{itemize}

\noindent
$S_2=|0000100000000000\rangle$.
\begin{itemize}
\item Compute then $S=S_1-S_2$, which here becomes
\end{itemize}

\noindent
\[
S=|0111000000000000\rangle
\]
\begin{itemize}
\item Perform then the logical AND operation of $S$ with the word containing 
\end{itemize}

\noindent
\[
\Phi_{0,3,6,10,13} = |1001001000100100\rangle,
\]
which results in $|0001000000000000\rangle$. Counting the number of $1-$bits gives the phase.  Here you need however an algorithm for bitcounting. Several efficient ones available. 

